\chapter{School Culture and Pupil Progress: Evidence from English Secondary Schools}
\label{ch:paper2}
\resetcounterformats

\begin{center}
  \textit{Damian Phelan%
    \thanks{Department of Economics, University College London.
      [\textit{INSERT: acknowledgements.}]}}
\end{center}

\begin{chapterabstract}
Does the culture of a school affect how much progress its pupils make? In
this paper I estimate the association between the warmth and strictness of
English secondary schools, measured by direct observation in the
\NTierOne{} schools I visited, and Progress~8, which is the Department for
Education's value-added measure of pupil progress. The regressions condition
on prior attainment, intake composition, school type, size, location and the
school's pre-pandemic inspection grade. There are three main results. The
first is that both dimensions predict progress, with warmth somewhat the
stronger of the two. A one standard deviation increase in observed warmth is
associated with Progress~8 that is $\hat{\beta}_W = \PrimarySpecW$ of a
grade higher, and a one standard deviation increase in observed strictness
with $\hat{\beta}_S = \PrimarySpecS$ of a grade ($n = \PrimarySpecN$). These
are significant at the one and five per cent levels respectively, and they
are stable across specifications. The second is that the association does
not run through the quality of teaching that was observed, which enters with
a coefficient of $0.015$ and leaves both culture coefficients unchanged. The
third is that the two dimensions add to one another rather than amplifying
each other: the schools that are both warm and strict sit $0.40$ of a grade
above the rest, and the interaction between the two is zero. Replacing the
observed scores with the headteacher's own account of the school removes the
warmth association, while the strictness score read from inspection reports
reproduces the strictness result across \NNationalStrictness{} schools with
the inspection grade held constant ($\BetaNationalStrictness$).
\end{chapterabstract}

% ---------------------------------------------------------------
\section{Introduction}
\label{sec:p2_intro}

Does the culture of a school affect how much progress its pupils make? In
this paper I estimate the association between the warmth and strictness of
English secondary schools, measured by direct observation, and the progress
their pupils make between the end of primary school and their GCSEs. I use
the enacted culture scores that were constructed in the previous chapter for
the \NTierOne{} visited schools, together with the Progress~8 measure
published by the Department for Education, which records each school's
average pupil progress relative to pupils nationally with the same prior
attainment.

The question belongs to a long tradition of asking whether the school a
child attends makes a difference beyond what the child brings to it, and if
so, through what. The tradition began with a negative answer.
\citet{coleman1966} found that school resources explained little of the
variation in attainment once family background was taken into account, and
\citet{hanushek1997}, reviewing the evidence three decades later, found no
consistent relationship between expenditure and performance. The class-size
experiments and natural experiments that followed
\citep{krueger1999,krueger2001,angrist1999} showed that resources can
matter, and English evidence finds small positive effects of spending that
are concentrated among disadvantaged pupils
\citep{levacic2005,jenkins2006,holmlund2010,machin2010}. However, these
effects are modest when set against the size of the school-to-school
variation that they leave unexplained. Two lines of work have located that
variation elsewhere. \citet{rutter1979}, studying twelve inner-London
secondary schools, placed it in what they called ethos, meaning the
expectations and routine practice that a school sustains, and the
effective-schools syntheses that followed listed an orderly environment and
a supportive climate among its correlates
\citep{edmonds1979,sammons1995,teddlie2000}. The teacher-quality literature
placed it in the classroom: \citet{rivkin2005} and \citet{chetty2014} show
that variation in teacher effectiveness is large and consequential, and
English estimates \citep{slater2012,murphy2011} are of the same order. What
neither line of work settles is whether a school's culture, as distinct from
the quality of its individual teachers, predicts how much progress its pupils
make. That is the question here.

Two strands of the literature are closest to this paper. The first is the
authoritative-school literature \citep{pellerin2005,gregory2010}, which
builds on the parenting typology of \citet{maccoby1983}, in which
responsiveness and demandingness each carry an independent effect on
children's outcomes \citep{majumder2015,zhang2020}. It finds that schools
which combine support with structure produce better behavioural and
engagement outcomes, but it has rarely had an academic-progress outcome or a
direct measure of both dimensions. The second is the literature on
high-expectation, high-structure charter and academy models, which finds
large effects of a bundle of practices. In the United States, effective
charter schools are marked by strict discipline and high expectations
\citep{abdulkadiroglu2011,angrist2010}; an index of five such practices
explains about half of the variation in charter effectiveness, while class
size and spending explain none of it \citep{dobbie2013}; and the bundle
transplanted into ordinary public schools raises mathematics attainment
\citep{fryer2014}. \citet{chabrier2016} and \citet{epple2016} review the
wider evidence, in which most charter schools show no average advantage and
the effective ones share this profile. In England, the first wave of
academies, which targeted under-performing schools and gave them autonomy,
raised attainment \citep{machin2011,eyles2016}, while the later converter
academies show much weaker effects \citep{eyles2019}. This suggests that
autonomy itself is not the operative channel. \citet{machin2020} show that
the early academies' higher exclusion rates reflected a genuine change in
discipline rather than the removal of weak pupils. In all of this work,
strict behaviour norms and warm, individual relationships sit together
inside one model, and the studies compare that model with its absence. This
paper measures the two dimensions separately, in ordinary schools rather
than in a selected model, and asks what each of them contributes.

The outcome is a value-added measure, and value-added measures have their
critics. \citet{gorard2010,gorard2013} argue that school-level estimates are
dominated by error and unstable from year to year. I use Progress~8
nonetheless, for three reasons. It is the measure that the Department for
Education publishes and by which schools are judged; it conditions on the
single strongest predictor of GCSE performance, which is prior attainment at
Key Stage~2; and it allows the visited schools to be placed against every
school in the country on the same footing. The concern about instability is
met by averaging two cohorts and by showing, in the appendix, that a
school's Progress~8 is in fact stable across years relative to the
differences between schools.

The design is an association, with several features that strengthen it
short of identification, and these are set out in full in the empirical
specification below. The outcome is value-added, so it already conditions on
where pupils started. The regressions condition on intake composition, prior
attainment, school type, size, location and the school's pre-pandemic
inspection grade. The sample excludes schools that admit pupils only at 13
or 14. Finally, the strictness result is reproduced across the national
population with a text-based measure. What the design cannot do is rule out
the possibility that some persistent characteristic of a school produces
both its culture and its results, and for that reason every estimate is
presented as an association.

There are three main results. First, both dimensions predict progress, with
warmth somewhat the stronger of the two. A one standard deviation increase
in observed warmth is associated with Progress~8 that is $0.15$ of a grade
higher, and a one standard deviation increase in observed strictness with
$0.11$ of a grade. These are significant at the one and five per cent levels
respectively, and both are stable across every treatment of the
inspection-grade control. Entered on its own, each dimension carries a
coefficient of about $0.2$, so the two share some variance but each carries
an independent signal. Second, the association does not run through the
quality of teaching that was observed. Adding the visit's teaching score
leaves both culture coefficients unchanged, and the teaching coefficient
itself is close to zero. Third, the schools that combine high warmth with
high strictness outperform the rest by two-fifths of a grade, while the
interaction between the two dimensions is zero. In other words, the
dimensions add to one another rather than amplifying each other. Replacing
the observed scores with the headteacher's own account of the school removes
the warmth association entirely, and the strictness score read from
inspection reports reproduces the strictness result across the national
population with the inspection grade held constant.

The rest of the paper is organised as follows. Section~\ref{sec:p2_data}
describes the outcome, the culture measures, the controls and the sample.
Section~\ref{sec:p2_spec} sets out the specification and what it can and
cannot identify. Section~\ref{sec:p2_results} reports the results and their
robustness. Section~\ref{sec:p2_conclusion} concludes.

% ============================================================
\section{Data}
\label{sec:p2_data}

\textbf{Outcomes.} The outcome is Progress~8, which is the Department for
Education's school-level value-added measure. For each pupil, attainment
across eight GCSE slots, weighted towards English and Maths, is compared with
the national average attainment of pupils who left primary school with the
same Key Stage~2 result, and the differences are then averaged across the
school's cohort. A school at zero is exactly average for its intake. A score
of $+0.5$ means that its pupils achieved half a grade more per subject than
pupils elsewhere with the same starting point, and $-0.5$ means half a grade
less. Nationally the measure has a standard deviation across schools of
about $0.5$ and an interquartile range of about two-thirds of a grade. This
means that the difference between a school at the first quartile and one at
the third is roughly two-thirds of a grade in every subject for every pupil.
I report overall Progress~8 and its English and Maths components, which are
the cleanest measures the tables offer because every pupil sits both of
those examinations. The remaining components, EBacc and Open, are partly
determined by which subjects a school enters its pupils for, and they are
reported in \cref{app:3A:tables} together with a direct analysis of subject
entry rates; \cref{app:3A:p8mechanics} sets out how the components are
constructed. The primary outcome is the two-year average of the 2022/23 and
2023/24 cohorts, which are the two post-pandemic years for which a Key
Stage~2 baseline exists. Averaging the two reduces the year-to-year noise in
a single cohort's figure. The 2024/25 cohort has no baseline, because its
Key Stage~2 tests were cancelled in 2020, and \cref{app:3A:tables} reports a
validated pseudo-progress measure for it.

\textbf{Culture measures.} The explanatory variables are the enacted warmth
and strictness scores of the previous chapter, that is, the visit-based
measures on a 0--10 scale for the \NTierOne{} visited schools. They are
entered per standard deviation so that the two coefficients are comparable.
Across the visited schools enacted warmth has a mean of $6.6$ and a standard
deviation of $1.0$, and enacted strictness a mean of $7.0$ and a standard
deviation of $0.9$, so a standard deviation is about one point on the
ten-point scale. Two properties of these scores, which were established in
the previous chapter, bear on how the estimates below should be read. The
first is that a school's score is an average over lessons that differ from
one another more than schools differ from each other (within-school lesson
variation is about three times the between-school variation for in-lesson
strictness). This means that ``a strict school'' here is a school whose
lessons are orderly on average, and not one in which every lesson is. The
second is that the scores carry measurement error, with a reliability of
about 0.7. This biases the coefficients on them towards zero rather than
away from it, so the estimates that follow are conservative. About a third
of the observed lessons were rated by one researcher rather than two, and the
robustness section checks that the estimates do not depend on those schools.
The visit's teaching-practice score, built from the ten lesson items on
questioning, feedback, explanation and structure, enters as a benchmark. It
is the most obvious rival explanation for any association between culture
and progress, since a school whose teachers are warm and orderly may also be
a school whose teachers teach well. The espoused scores, from the headteacher
interview, enter one robustness specification, and the inspection-report
strictness score enters the national extension.
\Cref{fig:ofsted_distributions} shows its distribution, and that of the
companion warmth score, across the national population.

\begin{figure}[htbp]
  \centering
  \includegraphics[width=\textwidth]{fig_A2_ofsted_distributions.pdf}
  \caption{Distributions of the inspection-report warmth (left) and strictness (right) scores
    across the national population (1--5 scale). Dashed lines mark means.
    The strictness score enters the national extension; the warmth score is
    not used as a predictor, for the reasons given in the previous chapter.}
  \label{fig:ofsted_distributions}
\end{figure}

\textbf{Controls.} All specifications condition on the same set of
controls. The first is the mean Key Stage~2 score of the cohort, which guards
against the possibility that Progress~8 is not fully neutral to intake. The
next are the shares of pupils eligible for free school meals, with English as
an additional language and with special educational needs, which are the
composition measures most strongly associated with progress. Then come log
school size; academy status, urban location and selective admissions; and
the number of years since the school's last inspection, since a school that
was inspected long ago is described by an older report. Finally, the
school's overall inspection grade as at August 2019 enters as a set of
dummies. I use the 2019 grade rather than the current one because a
contemporaneous grade is partly an outcome of the culture being measured. I
enter it at all because it is the best available summary of how the school
was regarded before the outcome cohorts arrived. Six visited schools
converted to academies after 2019 and carry no grade of their own. For five
of them the predecessor school's 2019 grade is used, and the sixth is
dropped where the control is required.

\textbf{Sample.} Two exclusions apply. The first is schools whose statutory
entry age is 13 or above (university technical colleges, studio schools and
similar; 98 schools nationally, two of them among those visited), which are
excluded from every outcome regression. Progress~8 measures progress from
age 11, and these schools enrol their pupils at 13 or 14, so their figure
attributes to them progress that was made elsewhere. The second is one
visited school that lacks the years-since-inspection control. The primary
estimation sample is therefore $n = \PrimarySpecN$. \Cref{tab:outcome_stats}
and \cref{tab:control_stats} report descriptive statistics for the outcomes
and the controls. The visited schools' mean Progress~8 is about a quarter of
a grade above zero, with a standard deviation a little below the national
one, which is the restricted range noted in the previous chapter. On the
controls the visited schools are close to the interviewed ones and to the
country as a whole: 29 per cent of pupils eligible for free school meals, 16
per cent with special educational needs, three-quarters academies, nine in
ten urban, one in twenty selective, and a 2019 grade distribution of roughly
a quarter Outstanding, five-eighths Good and one in nine below Good. The one
visible difference is the share of pupils with English as an additional
language, 28 per cent against 21 per cent in the interviewed sample, which
reflects the London weighting, and it is controlled for throughout. The
correlations between the culture scores and the controls are all below
$|0.3|$ (\cref{tab:score_controls_corr}). This means that the culture measures are not close proxies for intake or
school type, and that the coefficients below cannot simply be reflecting
differences in pupil composition.

\begin{table}[htbp]
  \centering
  \caption{Summary statistics: Progress~8 component outcomes by data tier}
  \label{tab:outcome_stats}
  \small
  \begin{tabular}{llrrrrr}
    \toprule
    Outcome & Tier & $N$ & Mean & SD & Min & Max \\
    \midrule
    \multirow{3}{*}{P8 English (2-yr avg)} & Tier~1: Visited (enacted) & 102 & 0.263 & 0.505 & -0.750 & 1.545 \\
     & Tier~2: Interview (303) & 299 & 0.092 & 0.488 & -1.485 & 1.545 \\
     & National & 3240 & -0.017 & 0.519 & -4.035 & 2.360 \\
    \addlinespace[3pt]
    \multirow{3}{*}{P8 Maths (2-yr avg)} & Tier~1: Visited (enacted) & 102 & 0.277 & 0.446 & -0.595 & 1.345 \\
     & Tier~2: Interview (303) & 299 & 0.092 & 0.473 & -1.370 & 1.345 \\
     & National & 3240 & -0.017 & 0.487 & -2.030 & 2.660 \\
    \addlinespace[3pt]
    \multirow{3}{*}{P8 EBacc (2-yr avg)} & Tier~1: Visited (enacted) & 102 & 0.307 & 0.543 & -0.790 & 1.490 \\
     & Tier~2: Interview (303) & 299 & 0.115 & 0.571 & -1.960 & 1.490 \\
     & National & 3240 & -0.020 & 0.587 & -4.950 & 2.655 \\
    \addlinespace[3pt]
    \multirow{3}{*}{P8 Open (2-yr avg)} & Tier~1: Visited (enacted) & 102 & 0.182 & 0.520 & -1.095 & 1.435 \\
     & Tier~2: Interview (303) & 299 & 0.054 & 0.556 & -3.095 & 1.500 \\
     & National & 3240 & -0.034 & 0.573 & -5.140 & 2.220 \\
    \addlinespace[3pt]
    \multirow{3}{*}{Att8 English (2024--25)} & Tier~1: Visited (enacted) & 102 & 10.685 & 1.744 & 8.000 & 17.300 \\
     & Tier~2: Interview (303) & 300 & 10.278 & 1.694 & 6.700 & 17.300 \\
     & National & 3274 & 10.049 & 1.702 & 0.100 & 17.300 \\
    \addlinespace[3pt]
    \multirow{3}{*}{Att8 Maths (2024--25)} & Tier~1: Visited (enacted) & 102 & 9.918 & 1.960 & 6.700 & 17.400 \\
     & Tier~2: Interview (303) & 300 & 9.510 & 1.899 & 5.200 & 17.400 \\
     & National & 3274 & 9.263 & 1.879 & 0.300 & 17.600 \\
    \addlinespace[3pt]
    \multirow{3}{*}{Att8 EBacc (2024--25)} & Tier~1: Visited (enacted) & 102 & 14.835 & 3.206 & 9.500 & 26.500 \\
     & Tier~2: Interview (303) & 300 & 14.157 & 3.117 & 6.200 & 26.500 \\
     & National & 3274 & 13.696 & 3.103 & 0.500 & 26.600 \\
    \addlinespace[3pt]
    \multirow{3}{*}{Att8 Open (2024--25)} & Tier~1: Visited (enacted) & 102 & 14.722 & 2.913 & 10.300 & 26.300 \\
     & Tier~2: Interview (303) & 300 & 14.189 & 2.860 & 4.900 & 26.300 \\
     & National & 3274 & 13.869 & 2.830 & 0.000 & 26.300 \\
    \bottomrule
  \end{tabular}
  \begin{minipage}{\linewidth}
    \vspace{2pt}
    \footnotesize
    \textit{Notes:} P8 outcomes are 2-year averages of 2022--23 and 2023--24 component scores where both
    years are available; single-year used where only one is present. Att8 2024--25 components are the
    contemporaneous robustness outcome; no Progress~8 is available for this cohort because the cohort were
    in Year~6 in 2019--20 when Key Stage~2 assessments were cancelled. Tier~2 is the full interviewed
    sample and therefore contains Tier~1; both are subsets of the national sample, hence the different
    $N$ values across tiers.
  \end{minipage}
\end{table}


\begin{table}[htbp]
  \centering
  \caption{Summary statistics: control variables by data tier}
  \label{tab:control_stats}
  \small
  \begin{tabular}{lrrrrrr}
    \toprule
    & \multicolumn{3}{c}{Tier~1 ($N=102$, visited)} & \multicolumn{3}{c}{Tier~2 ($N=303$, interview)} \\
    \cmidrule(lr){2-4}\cmidrule(lr){5-7}
    Variable & $N$ & Mean & SD & $N$ & Mean & SD \\
    \midrule
    \multicolumn{7}{l}{\textit{Pupil composition and prior attainment}} \\
    \quad Mean KS2 score & 102 & 104.917 & 2.909 & 300 & 104.627 & 2.714 \\
    \quad FSM eligible (\%) & 102 & 29.2 & 15.7 & 303 & 28.0 & 14.7 \\
    \quad EAL share (\%) & 102 & 27.8 & 26.2 & 300 & 21.3 & 22.2 \\
    \quad SEN share (\%) & 102 & 15.5 & 7.3 & 300 & 15.6 & 7.3 \\
    \quad Log school size & 102 & 6.964 & 0.368 & 303 & 6.938 & 0.398 \\
    \quad Yrs since Ofsted & 101 & 4.797 & 3.891 & 295 & 4.919 & 4.006 \\
    \addlinespace[3pt]
    \multicolumn{7}{l}{\textit{School characteristics (share)}} \\
    \quad Academy school & 102 & \multicolumn{2}{c}{73.5\%} & 303 & \multicolumn{2}{c}{71.6\%} \\
    \quad Urban location & 102 & \multicolumn{2}{c}{90.2\%} & 303 & \multicolumn{2}{c}{88.4\%} \\
    \quad Selective admissions & 102 & \multicolumn{2}{c}{93.1\%} & 303 & \multicolumn{2}{c}{92.4\%} \\
    \addlinespace[3pt]
    \multicolumn{7}{l}{\textit{Pre-COVID Ofsted grade (as at 31~August~2019)}} \\
    \quad Grade~1 (Outstanding) & 25 & \multicolumn{2}{c}{26.3\%} & 63 & \multicolumn{2}{c}{23.6\%} \\
    \quad Grade~2 (Good) & 59 & \multicolumn{2}{c}{62.1\%} & 170 & \multicolumn{2}{c}{63.7\%} \\
    \quad Grade~3 (Requires improvement) & 10 & \multicolumn{2}{c}{10.5\%} & 27 & \multicolumn{2}{c}{10.1\%} \\
    \quad Grade~4 (Inadequate) & 1 & \multicolumn{2}{c}{1.1\%} & 7 & \multicolumn{2}{c}{2.6\%} \\
    \quad No pre-COVID grade & 7 & & & 36 & & \\
    \bottomrule
  \end{tabular}
  \begin{minipage}{\linewidth}
    \vspace{2pt}
    \footnotesize
    \textit{Notes:} KS2 is the mean scaled score of the 2023--24 KS4 cohort at Key Stage~2.
    FSM is the share eligible for the pupil premium (FSM6 definition); EAL and SEN are KS4
    cohort shares from the 2023--24 performance tables. Pre-COVID Ofsted grade is the overall
    effectiveness grade from the Ofsted Management Information snapshot as at 31~August~2019,
    used as the primary endogeneity-safe control; the contemporary 2024 grade enters sensitivity
    analyses only. Tier~1 (102 schools) have both visit and interview data; Tier~2 (303 schools)
    includes Tier~1.
  \end{minipage}
\end{table}

As the previous chapter records, the visited schools lean towards stronger
schools than the national population, so the estimates apply most directly
to a population somewhat above the national average, and the restricted
range of both culture and outcomes works against finding associations rather
than for them.

% ============================================================
\section{Empirical Specification}
\label{sec:p2_spec}

The primary specification is
%
\begin{equation}
  \label{eq:main}
  Y_s = \alpha + \beta_W W_s + \beta_S S_s + \mathbf{X}_s'\gamma
        + \sum_g \delta_g \,\mathbb{1}[G_s = g] + \varepsilon_s ,
\end{equation}
%
where $Y_s$ is the two-year average Progress~8 (overall, English or Maths)
of school $s$, $W_s$ and $S_s$ are its enacted warmth and strictness scores
in standard-deviation units, $\mathbf{X}_s$ is the set of controls listed
above and $G_s$ is its 2019 inspection grade. Standard errors are
heteroskedasticity-consistent (HC3), which is the small-sample variant that
is appropriate when $n$ is near 100. I estimate three stages:
equation~(\ref{eq:main}); the same equation with the teaching-practice score
$T_s$ in place of $W_s$ and $S_s$; and the equation with all three. Each
dimension is also entered on its own. The national extension replaces $W_s$
and $S_s$ with the inspection-report strictness score, standardised, on all
schools with an outcome, and is estimated both with and without the grade
dummies.

The estimates are associations, and it is worth being exact about what the
design does and does not achieve. Because the outcome is value-added, the
comparison is between schools whose pupils started from the same place. This
means that the first and largest confound in any school-effects estimate,
namely that stronger schools take stronger pupils, is removed by
construction, and the Key Stage~2 control guards against what remains.
Because the culture measures come from visits that were mostly made in the
year after the outcome cohorts sat their examinations, reverse causation in
the ordinary sense, that good results produce a warm and orderly school, is
not ruled out by the timing. What the design relies on instead is that culture is persistent, so that
what was observed in 2024/25 is a reasonable measure of what the 2022/23
and 2023/24 cohorts experienced. The appendix
supports that reading in two ways. A school's Progress~8, and the structural
characteristics that shape it, are stable across years relative to the
differences between schools, and the one cohort that was measured at the
same time as the fieldwork gives coefficients within $0.05$ of the primary
ones. Because the 2019 grade is held constant, the comparison is between
schools that the inspectorate regarded alike before the outcome cohorts
arrived, so a school's reputation, and whatever that reputation draws in
terms of staff and families, is partly absorbed. Finally, because the sample excludes schools that enrol at 13 or 14, no
school is credited with progress that was made at another school.

What the design cannot do is rule out the possibility that some persistent
characteristic of a school produces both its culture and its results. This
might be a headteacher who is effective in ways that the visit does not
measure, or a stable school reputation that draws in both committed staff
and motivated families. Two features of the data bear on this without
resolving it. The teaching-practice score is observed alongside the culture
scores, so the most obvious such characteristic, that warm and orderly
schools are simply schools whose teachers teach well, can be tested directly
rather than assumed away. Secondly, the strictness result reproduces across the
national population using a text-based measure with the inspection grade
held constant, so it is not an artefact of which schools agreed to be
visited. The sorting concern that remains, which is that strict schools may
retain fewer pupils with social, emotional and mental-health needs and gain
in measured progress by that route, is tested by adding a pre-period measure
of that share. The estimates are nonetheless presented as associations
throughout, and no coefficient is described as an effect.

% ============================================================
\section{Results}
\label{sec:p2_results}

\textbf{Culture predicts progress.} \Cref{tab:spec_ladder} reports the
primary estimates. In the primary specification observed warmth has a
coefficient of $\hat{\beta}_W = \PrimarySpecW$ ($p = 0.002$) and observed
strictness a coefficient of $\hat{\beta}_S = \PrimarySpecS$ ($p = 0.007$) on
overall Progress~8, with $R^2 = 0.62$. The table shows the same equation
under four treatments of the grade control, from requiring an original 2019
grade ($n = 95$) to dropping the control altogether ($n = 100$), and the
coefficients move by less than $0.02$ across the four. Both dimensions
contribute, as the authoritative-school framework predicts, with warmth
somewhat the stronger of the two.

\begin{table}[htbp]\centering
\def\sym#1{\ifmmode^{#1}\else\(^{#1}\)\fi}
\small
\caption{The headline gold-standard model under four treatments of the pre-COVID
inspection-grade control. Each row is the same regression of average Progress 8 on
enacted warmth ($W$) and enacted strictness ($S$) with the full control set; the rows
differ only in how schools with a missing 2019 grade are handled. Schools with a
statutory admission age of 13 or above are excluded throughout.}
\label{tab:spec_ladder}
\begin{tabular}{lccc}
\toprule
Grade-control specification & $N$ & Warmth ($W$) & Strictness ($S$) \\
\midrule
Grade as recorded (missing dropped)      &  95 & 0.155\sym{***}  & 0.128\sym{***}  \\
                                         &     & (0.045)        & (0.047)        \\
\addlinespace
Predecessor-filled grade (primary)       &  99 & 0.143\sym{***}  & 0.131\sym{***}  \\
                                         &     & (0.045)        & (0.047)        \\
\addlinespace
Missing-grade category retained          & 100 & 0.143\sym{***}  & 0.131\sym{***}  \\
                                         &     & (0.045)        & (0.047)        \\
\addlinespace
No grade control                         & 100 & 0.140\sym{***}  & 0.127\sym{***}  \\
                                         &     & (0.046)        & (0.046)        \\
\bottomrule
\end{tabular}
\begin{minipage}{\linewidth}
\smallskip
\footnotesize\textit{Notes:} Standard errors in parentheses.
\sym{*} \(p<0.10\), \sym{**} \(p<0.05\), \sym{***} \(p<0.01\).
Outcome is average Progress 8; predictors are standardised. Seven visited academies
lacked a 2019 grade under their current URN; the primary specification fills five of
them from the predecessor school's grade. Including the late-entry schools leaves the
primary estimates essentially unchanged ($W = 0.139$, $p = 0.003$;
$S = 0.133$, $p = 0.006$; $N = 100$).
\end{minipage}
\end{table}


It is worth putting these magnitudes in context. A standard deviation of
enacted warmth is about one point on the ten-point scale, or half a point on
the five-point scale of every warmth item in every lesson observed. The
associated $0.15$ of a grade is about a quarter of a standard deviation of
school Progress~8 nationally. It is also about a quarter of the gap in mean
Progress~8 between schools graded Outstanding and Good in 2019 ($0.53$), and
about a quarter of the gap between the least and the most disadvantaged
quarter of schools by free-school-meal share ($0.62$). Two standard deviations on both dimensions together, which is roughly the
distance within the visited range from a school that is low on both
dimensions to one that is high on both, is associated with about half a
grade per subject per pupil. That is the size
of the difference that the performance tables treat as separating a
well-above-average school from an average one.

Entered on its own, each dimension has a coefficient of about $0.2$
(\cref{tab:univariate_ws}): $0.219$ for warmth and $0.199$ for strictness.
Entered jointly, warmth keeps two-thirds of its solo coefficient and
strictness about half. The two dimensions therefore share some variance, as
their correlation of $\CorrWarmthStrictnessEnacted$ implies, but each
contains substantial information that the other does not. This is the
empirical basis for treating them as two constructs. If there were a single
scale running from warm to strict, one coefficient would absorb the other.

\begin{table}[htbp]\centering
\def\sym#1{\ifmmode^{#1}\else\(^{#1}\)\fi}
\small
\caption{Warmth and strictness entered alone and together. Each column is a Progress 8
outcome; the first two rows enter one culture dimension at a time and the final two
rows enter both jointly, always with the full control set and the predecessor-filled
grade control. The drop from the univariate to the joint coefficients shows the two
dimensions share some variance, but each keeps an independent association when the
other is held fixed.}
\label{tab:univariate_ws}
\begin{tabular}{llccc}
\toprule
Specification & Term & Overall & English & Maths \\
\midrule
$S$ only   & Strictness ($S$) & 0.195\sym{***} & 0.170\sym{***} & 0.156\sym{***} \\
           &                  & (0.033)        & (0.039)        & (0.033)        \\
\addlinespace
$W$ only   & Warmth ($W$)     & 0.206\sym{***} & 0.207\sym{***} & 0.164\sym{***} \\
           &                  & (0.035)        & (0.037)        & (0.039)        \\
\midrule
Joint      & Warmth ($W$)     & 0.132\sym{***} & 0.158\sym{***} & 0.104\sym{**}  \\
           &                  & (0.042)        & (0.041)        & (0.050)        \\
\addlinespace
Joint      & Strictness ($S$) & 0.117\sym{***} & 0.077\sym{*}   & 0.095\sym{**}  \\
           &                  & (0.039)        & (0.042)        & (0.042)        \\
\midrule
$N$        &                  & 100            & 100            & 100            \\
\bottomrule
\end{tabular}
\begin{minipage}{\linewidth}
\smallskip
\footnotesize\textit{Notes:} Standard errors in parentheses.
\sym{*} \(p<0.10\), \sym{**} \(p<0.05\), \sym{***} \(p<0.01\).
Gold-standard visited tier, late-entry schools excluded. EBacc and Open components
are reported in the chapter appendix.
\end{minipage}
\end{table}


At the level of the components, warmth is significant for both subjects (English
$0.162$, $p < 0.001$; Maths $0.135$, $p = 0.011$), while strictness falls
short of significance in each taken alone (English $0.071$, $p = 0.11$;
Maths $0.076$, $p = 0.09$), so that its association with overall progress
is the sum of two parts with the same sign. Estimates for the two remaining
components, and the analysis of subject entry that they require, are in
\cref{app:3A:tables}.

\textbf{Teaching quality does not account for it.} \Cref{tab:stages23_trio}
reports the second and third stages. On its own, the visit's
teaching-practice score predicts overall Progress~8 at $0.206$ ($p <
0.001$), but with a lower $R^2$ ($0.55$) than the two culture scores. When it
is added to the primary specification it enters at $0.015$ ($p = 0.82$), and
the culture coefficients are unchanged ($0.137$ and $0.103$).

There are two possible readings of this result, and the data do not allow
a choice between them. The first reading is that culture and instruction
are distinct predictors of progress. On this reading, what a school is like
between and around lessons, and how orderly its lessons are, matter over and
above how well the lessons are taught, and the culture association is not
simply the teaching-quality association under another name. The second
reading is that the teaching score is the noisiest of the three visit
measures, because it is built from ten items on which the observers agree
least, and that its signal is absorbed by the two culture scores with which
it is correlated. On this reading the null result on teaching understates
what instruction contributes. Three facts favour the first reading, although
they do not settle the matter. Firstly, the teaching score on its own
predicts progress well, so it is clearly not just noise. Secondly, the
culture coefficients do not move at all when it enters, whereas if
absorption were taking place one would expect to see movement in both
directions. Thirdly, the culture association runs partly through
out-of-lesson warmth, which no measure of classroom teaching could pick up.
What would separate the two readings is a measure of teaching quality that
is independent of the visit, and no such measure exists for these schools.
The claim I make is therefore the narrower one, which is that conditioning
on observed teaching quality does not account for the culture association.

\begin{table}[htbp]
\centering
\small
\caption{Stages 2 and 3 under the primary specification: teaching quality as
the sole culture predictor (Stage~2), and teaching quality added to the
warmth-and-strictness model (Stage~3). Coefficients are per standard
deviation of the predictor; HC3 robust standard errors; $n = 100$. The full
five-outcome versions, estimated on the previous specification, are reported
in the appendix.}
\label{tab:stages23_trio}
\begin{tabular}{lccc}
\toprule
 & \textit{Overall P8} & \textit{English} & \textit{Maths} \\
\midrule
\multicolumn{4}{l}{\textit{Stage 2: teaching quality only}} \\
\addlinespace[2pt]
Teaching $T$ & 0.202\sym{***} & 0.184\sym{***} & 0.159\sym{***} \\
             & (0.000) & (0.000) & (0.001) \\
$R^2$        & 0.562 & 0.560 & 0.484 \\
\addlinespace[6pt]
\multicolumn{4}{l}{\textit{Stage 3: warmth, strictness and teaching}} \\
\addlinespace[2pt]
Warmth $W$   & 0.121\sym{*}  & 0.165\sym{**} & 0.098 \\
             & (0.033) & (0.006) & (0.108) \\
Strictness $S$ & 0.113\sym{**} & 0.080 & 0.093\sym{*} \\
             & (0.005) & (0.064) & (0.047) \\
Teaching $T$ & 0.020 & $-$0.010 & 0.011 \\
             & (0.758) & (0.879) & (0.872) \\
\bottomrule
\end{tabular}
\begin{minipage}{0.8\linewidth}
\vspace{4pt}\scriptsize
\textit{Notes}: $p$-values in parentheses. \sym{*} $p<0.05$, \sym{**}
$p<0.01$, \sym{***} $p<0.001$. Primary specification throughout: full
control set, predecessor-filled pre-COVID grade, late-entry schools
excluded.
\end{minipage}
\end{table}


\textbf{The dimensions add.} \Cref{fig:quadrant_ws} plots the visited
schools by their two culture scores, shaded by progress. Schools in the warm
and strict quadrant are predominantly above-average performers, and the
other three quadrants are all populated, which is what makes it possible to
test the combination. Formally, the warm-and-strict quadrant sits $0.40$ of
a grade above the rest of the sample after controls ($p < 0.001$), while the
interaction between the two standardised scores is zero ($0.038$, $p =
0.27$; \cref{tab:robustness_overall}). In other words, warmth and strictness
add to one another rather than amplifying each other. The best schools have
both because each of them helps, and not because of any combination effect,
and a school that is strong on one dimension still gains from the other. Breaking
the visit scores down into their five observational blocks
(\cref{tab:subscores}) shows that the association runs through warmth
outside lessons and orderly conduct within them. Out-of-lesson warmth
($0.111$, $p = 0.006$) and in-lesson strictness ($0.115$, $p = 0.008$)
survive joint entry, and their counterparts do not. This means that the
relational quality that goes with progress is the one visible in corridors
and at break, and the order that goes with it is the one visible in lessons.

\begin{figure}[htbp]
  \centering
  \includegraphics[width=0.85\textwidth]{fig_quadrant_ws.pdf}
  \caption{Enacted warmth and strictness for the visited schools, shaded by
  Progress~8 (two-year average). Dashed lines mark sample medians. The
  adjusted quadrant comparison is in \cref{tab:typology}.}
  \label{fig:quadrant_ws}
\end{figure}

\textbf{Espoused culture does not predict progress.} Substituting the
headteacher's espoused scores for the observed ones on the same schools
reduces the warmth coefficient to $0.051$ ($p = 0.29$), while espoused
strictness retains an association of the same size as the observed one
($0.120$, $p = 0.018$). The asymmetry found in the previous chapter's measurement appears again
here in prediction. What a headteacher says about warmth contains no
information about progress, and what they say about behaviour systems
contains some. This is the pattern that the previous chapter's
finding predicts, since espoused warmth was shown there to be only weakly
related to observed warmth, and espoused strictness more closely related to
observed strictness. It is also a check on the observed measures as much as
on the espoused ones. Had the espoused scores predicted progress as well as
the observed ones, the visits would have added nothing that an hour with the
headteacher could not supply. Two related results are in
\cref{app:3A:tables}. The gap between a school's espoused and enacted scores
predicts nothing beyond the enacted level, and the association is not
confined to the visited schools' own report on themselves, since parents'
survey responses on their children's happiness and safety are positively
associated with progress ($0.14$ per standard deviation), with the caveat
that satisfaction and results are jointly determined.

\textbf{National extension.} The inspection-report strictness score, which
the previous chapter showed to track observed strictness at $r =
\CorrOfstedStrictness$, is available for \NNationalStrictness{} schools with
an outcome, which is some thirty times the size of the visited sample. On
the full control set without the grade it has a coefficient of
$\BetaNationalStrictness$ ($p < 0.001$, $R^2 = 0.52$;
\cref{tab:national_strictness}). With the 2019 grade added ($n = 2{,}746$)
the coefficient is $0.140$, and no component estimate moves by more than
$0.005$. Two things follow from this. The first concerns
the score itself. A score read from an inspection report might be thought to be a proxy for
the inspectorate's overall approval, in which case the association would
really be the grade's association. However, the coefficient is unchanged by
conditioning on the grade, and a proxy for the grade would lose most of its
coefficient when the grade entered. The score therefore captures what the
report says about conduct and enforcement over and above the verdict, and
it is this that predicts progress. The second concerns the visited
sample. The visited schools agreed to be visited, and the previous chapter
showed that they lean towards stronger schools. The national coefficient is
estimated on every school with a report and an outcome, and it is of the
same order as the visited strictness coefficient in standard-deviation
units, if anything larger. The strictness result is therefore not a property
of the schools that agreed to take part. What the extension does not provide
is the same reassurance for warmth, since no text source contains observed
warmth, and nor does it provide identification, since a text-derived measure
is open to the same omitted-variable concern as an observed one. Two further text-derived
results are consistent with the previous chapter's measurement findings and
are reported in \cref{tab:llm_p8_matrix}: the behaviour-policy warmth score
is significantly \emph{negatively} associated with progress, and the website
warmth and strictness scores are not associated with it at all.

\begin{table}[htbp]\centering
\def\sym#1{\ifmmode^{#1}\else\(^{#1}\)\fi}
\caption{National extension: Ofsted LLM strictness and Progress~8}
\label{tab:national_strictness}
\resizebox{\textwidth}{!}{%
\begin{tabular}{l*{5}{c}}
\toprule
 & Overall & English & Maths & EBaC & Open \\
\midrule
\multicolumn{6}{l}{\textit{Panel A: published specification (Ofsted grade omitted)}} \\
Strictness (Ofsted LLM, 1--5 scale) & 0.166\sym{***} & 0.152\sym{***} & 0.142\sym{***} & 0.171\sym{***} & 0.190\sym{***} \\
 & (0.008) & (0.009) & (0.008) & (0.010) & (0.010) \\
$N$ & 3{,}147 & 3{,}147 & 3{,}147 & 3{,}147 & 3{,}147 \\
$R^2$ & 0.522 & 0.455 & 0.471 & 0.499 & 0.423 \\
\addlinespace[6pt]
\multicolumn{6}{l}{\textit{Panel B: adding the pre-COVID (2019) Ofsted grade}} \\
Strictness (Ofsted LLM, 1--5 scale) & 0.167\sym{***} & 0.154\sym{***} & 0.143\sym{***} & 0.174\sym{***} & 0.190\sym{***} \\
 & (0.009) & (0.009) & (0.009) & (0.011) & (0.010) \\
$N$ & 2{,}823 & 2{,}823 & 2{,}823 & 2{,}823 & 2{,}823 \\
$R^2$ & 0.534 & 0.463 & 0.477 & 0.510 & 0.435 \\
\bottomrule
\end{tabular}%
}
\begin{minipage}{\linewidth}
  \vspace{4pt}
  \footnotesize
  \textit{Notes:} Heteroskedasticity-robust (HC3) standard errors in
  parentheses. \sym{*} \(p<0.10\), \sym{**} \(p<0.05\), \sym{***} \(p<0.01\).
  All specifications control for prior attainment (KS2), FSM, EAL, SEN, log
  cohort size, years since inspection, academy status, urban location and
  selective status. Panel~A is the specification reported in the text, which
  omits the Ofsted grade because the strictness score is read from the same
  inspection report. Panel~A includes 3{,}147 schools. Panel~B adds the pre-COVID
  (2019) overall Ofsted grade as a set of dummies---the same control used in
  the primary Tier~1 specification---and is estimated on the 2{,}823 schools with
  a 2019 grade on record, 324 fewer than Panel~A. The strictness
  coefficient is materially unchanged throughout, indicating that it is not
  simply recording Ofsted's overall approval of the school.
\end{minipage}
\end{table}


\textbf{Robustness.} \Cref{tab:robustness_overall} reports the primary
estimate under alternative specifications, all on the primary sample, and
each of them addresses a particular doubt. The grade control might be doing
the work, or absorbing part of it. Dropping it leaves $0.153$ and $0.097$
($n = 100$). Averaging two cohorts might smooth over a result that holds in
only one of them. The single 2023/24 outcome year gives $0.155$ and $0.105$.
Strict schools might have fewer pupils with social, emotional and
mental-health needs, and gain in measured progress by that route. Adding the
2015/16 baseline share of such pupils, which was measured before any of the
culture was observed, gives $0.149$ and $0.090$ ($n = 94$), and the sorting
channel accounts for under a tenth of the strictness association
(\cref{tab:semh_mechanism}). Progress~8 might be peculiar. On the 2024/25
cohort's Attainment~8, which does not condition on intake, both dimensions
remain significant, both on the total of the four buckets ($1.40$, $1.32$
points) and on the English bucket alone ($0.35$, $0.16$, the last at $p =
0.12$).
The enacted scores might reflect something about the visits themselves. About
a third of the observed lessons were rated by one researcher rather than
two, those lessons were scored a little lower, and the share of such lessons
in a school's visit correlates $-0.25$ with its warmth score. Adding that
share as a control moves nothing ($0.150$ and $0.105$; the share itself
has a coefficient of $0.03$, $p = 0.76$). Restricting the sample to the 64 schools in
which every lesson was rated by two researchers gives $0.189$ for warmth
($p = 0.001$) and $0.071$ for strictness ($p = 0.15$). The warmth estimate is
larger on that subset and the strictness estimate smaller, with a standard
error a third larger on two-thirds of the sample, so the strictness result
is the one that depends on the full sample. All but the English-bucket and
the double-rated strictness cells leave both coefficients significant at the
five per cent level or better.

Two further checks concern timing and the headteacher. Within-school
Progress~8 is stable across years (the mean within-school standard deviation
is $0.20$, against a between-school standard deviation of about $0.5$),
which supports using culture measured after the outcome window as a proxy
for culture during it (\cref{tab:stability_p8}), and the 2024/25
pseudo-progress measure, which is the one cohort measured at the same time
as the fieldwork, gives coefficients within $0.05$ of their real-Progress~8
counterparts (\cref{tab:p8_proxy}). Restricting the sample to the 22 schools
whose headteacher is confirmed as unchanged since the last inspection, so
that the culture observed is the culture the outcome cohorts had, leaves the
strictness point estimates at or above the full-sample values and the warmth
estimates imprecise. That subsample is too small to be informative either
way (\cref{app:3A:tables}).

\begin{table}[htbp]\centering
\def\sym#1{\ifmmode^{#1}\else\(^{#1}\)\fi}
\caption{Robustness: Overall Progress 8}
\label{tab:robustness_overall}
\resizebox{\textwidth}{!}{%
\begin{tabular}{l*{6}{c}}
\toprule
                    &\multicolumn{1}{c}{(1)}&\multicolumn{1}{c}{(2)}&\multicolumn{1}{c}{(3)}&\multicolumn{1}{c}{(4)}&\multicolumn{1}{c}{(5)}&\multicolumn{1}{c}{(6)}\\
                    &\multicolumn{1}{c}{Primary}&\multicolumn{1}{c}{No grade}&\multicolumn{1}{c}{2023-24}&\multicolumn{1}{c}{Att8 2425}&\multicolumn{1}{c}{W$\times$S}&\multicolumn{1}{c}{SEMH ctrl}\\
\midrule
$W$                 &       0.127\sym{***}&       0.142\sym{***}&       0.138\sym{***}&       0.895\sym{*}  &      -0.106         &       0.130\sym{***}\\
                    &     (0.045)         &     (0.045)         &     (0.049)         &     (0.502)         &     (0.294)         &     (0.044)         \\
\addlinespace
$S$                 &       0.120\sym{**} &       0.099\sym{**} &       0.119\sym{**} &       1.767\sym{***}&      -0.099         &       0.114\sym{***}\\
                    &     (0.047)         &     (0.044)         &     (0.047)         &     (0.476)         &     (0.270)         &     (0.041)         \\
\midrule
N                   &          95         &         101         &          95         &          95         &          95         &          90         \\
r2                  &       0.595         &       0.568         &       0.601         &       0.883         &       0.599         &       0.638         \\
\bottomrule
\multicolumn{7}{l}{\footnotesize Standard errors in parentheses}\\
\multicolumn{7}{l}{\footnotesize \sym{*} \(p<0.10\), \sym{**} \(p<0.05\), \sym{***} \(p<0.01\)}\\
\end{tabular}%
}
\end{table}


\textbf{Implications.} Two implications bear on the current policy debate,
and they follow from the pattern of results rather than from any single
coefficient. The behaviour debate treats discipline as the part of school improvement
that does the work, and relational culture as secondary. These estimates
give relational culture at least the weight of discipline, they show no
trade-off between the two, and they place the warm-and-strict schools
two-fifths of a grade ahead. On these estimates, a school that tightens its
behaviour systems but lets relationships slip loses about as much as it
gains.
The second implication is that the association runs through relational
quality outside lessons and orderly conduct within them. This places the relational work that matters outside the classroom, in
corridors, at break and at the gate, which is where schools currently direct
the least observation and training effort. Both implications are conditional on the association being at
least partly causal, which the design supports but does not establish.

% ============================================================
\section{Conclusion}
\label{sec:p2_conclusion}

Observed school culture predicts pupil progress. Across the \NTierOne{}
English secondary schools that were visited, a standard deviation of
observed warmth is associated with Progress~8 that is $0.15$ of a grade
higher, and a standard deviation of observed strictness with $0.11$ of a
grade. These are significant at the one and five per cent levels
respectively, and both are stable across every reasonable specification.
The two dimensions contribute additively, with warmth somewhat the stronger,
and the schools that combine them outperform the rest by two-fifths of a
grade. Observed teaching quality does not account for the association. The
headteacher's own account of the school predicts nothing on warmth, and the
strictness score read from inspection reports reproduces the strictness
result nationally with the inspection grade held constant.

Against the literature the paper sits in, the results say three things. To
the school-effects tradition, they give a quantitative reading of two of the
correlates of ethos that Rutter and the effective-schools syntheses named,
measured by observation across a national population rather than in a
case-study sample. To the charter and academy literature, they say that the
two halves of the practice bundle, structure and support, each matter
when they are measured apart, and that neither is the whole of the effect.
To the authoritative-school literature, they supply the
academic-progress outcome that it has mostly lacked, and they find what the
typology predicts: both dimensions matter, and the schools with both do best.

The estimates are associations. They are strengthened by a value-added
outcome, by the ordering in time, by a rich set of controls and by their
stability across specifications, but they are not identified: a persistent
characteristic that produces both a school's culture and its results cannot
be ruled out with cross-sectional data. The visited sample is small and leans towards stronger schools, culture
was measured after the outcome window for most schools, and a single visit
measures a school with considerable error. A longitudinal design, with
culture observed before the outcome cohort arrived and again afterwards,
would provide evidence on the direction of the association, and a design
that observed culture changing within schools, as headteachers changed or
behaviour systems were introduced, would provide evidence on its cause. Within the
present limits, the results identify the relational and disciplinary
features of a school that go together with its pupils' progress, and they
show that one of the two, strictness, can be read at national scale from the
public record.

% ===============================================================
%  APPENDIX 3A
% ===============================================================
\clearpage
\startchapterappendix{3}
\begin{center}
  {\LARGE\bfseries Appendix to Chapter 3}\\[6pt]
  \rule{0.5\linewidth}{0.4pt}
\end{center}
\medskip
\noindent This appendix collects the supporting material for the analyses
reported in the chapter. It explains how the Progress~8 measures and the
subject entry rates are calculated, sets out the evidence on the stability of
school characteristics that underpins the choice of outcome years, and gives
the full five-outcome versions of the main results tables (the chapter body
reports overall, English and Maths Progress~8, and the EBaC and Open
components appear here) together with the secondary and exploratory
analyses.
\medskip

% - - - - - - - - - - - - - - - - - - - - - - - - - - - - - - -
\section{How the Progress 8 Measures Are Calculated}
\label{app:3A:p8mechanics}

Progress~8 measures the progress a pupil makes between the end of primary
school and GCSE, relative to pupils nationally who started from the same
point. Each pupil's GCSE results are counted across eight slots, their
attainment is compared with the national average attainment of pupils with
the same Key Stage~2 starting point, and the difference is averaged across
the school. A school at zero is exactly average; a school at $+0.5$ achieves
half a grade more per subject per pupil than would be expected from its
intake.

The eight slots are structured, and this structure is why the thesis
treats the components differently. English and Maths each fill two slots and
are double-weighted. Every pupil sits both, so these components measure
progress and nothing else. The EBaC component covers three slots that must be filled
from a restricted academic list: sciences, computer science, history,
geography, and modern foreign languages. The Open component covers the final
three slots, which may be filled by any approved qualification, including
arts and vocational subjects. For these last six slots the school makes
entry decisions, and an unfilled slot contributes zero progress to the
school's average. Consequently, a school's EBaC and Open components mix two
different things: how much progress pupils make in the subjects they enter,
and which subjects the school enters them for. A school that enters few
pupils for languages will carry empty EBaC slots and a depressed EBaC
component even if the pupils it does enter progress well.

This structure generates a theoretical prediction about warmth and
strictness. Schools higher on strictness, characterised by structured
behaviour management, high academic expectations, and intensive pupil
steering, are more likely to direct pupils into the EBaC pathway, and could
therefore raise their EBaC component through the entry-rate channel as well
as through progress itself. Schools higher on warmth might generate stronger
progress in the more elective Open component. The five-outcome versions of the main results tables in this appendix
provide evidence on that prediction, and the entry-rate analysis below
addresses the pathway channel directly.

The entry rates themselves are calculated from the pupil count columns of
the 2023--24 performance tables, because the Department for Education
stopped publishing entry rates in percentage form in recent releases. The
EBaC entry rate is the number of pupils entered for the full EBaC suite
divided by the number of pupils at the end of Key Stage~4, and the
humanities and languages rates are constructed in the same way from their
respective counts. Entry rates are available only at the subject-group
level; individual subjects within a group cannot be distinguished in the
data collected for this project.

% - - - - - - - - - - - - - - - - - - - - - - - - - - - - - - -
\section{Structural Stability of School Characteristics}
\label{app:3A:stability}

The intra-class correlation (ICC) decomposes total variance in a panel variable
into a between-school component and a within-school over-time component. An ICC
close to one indicates that schools are much more different from each other
than they are from their own past selves. In other words, year-on-year
fluctuation within a school is small relative to the persistent differences
between schools. ICC is estimated by
one-way random-effects ANOVA on the school-year panel, retaining only schools
with at least three years of non-missing data for the variable in question.

\Cref{tab:structural_stability} reports ICC estimates and mean absolute
year-on-year changes for nine panel variables covering pupil intake composition,
special educational needs prevalence, prior attainment, pupil progress, and
per-pupil expenditure. Six of nine variables exceed ICC~$= 0.69$; all
exceed~$0.49$. Intake characteristics (the share of pupils with English as
an Additional Language, school roll, and the mean KS2 prior attainment of the
entering cohort) are especially stable, with ICC values between 0.81 and 0.92
and small mean absolute year-on-year changes relative to their between-school
spread.
Special educational needs prevalence (EHC plan share, ICC~$= 0.67$; SEN support
share, ICC~$= 0.65$) and school-level Progress~8 (ICC~$= 0.73$) show moderate to
high stability. Per-pupil expenditure, at ICC~$= 0.52$, is the most volatile, which
reflects its sensitivity to changes in national funding policy. Even so, the
majority of the variation in expenditure is between schools rather than
within schools over time.

These patterns are consistent with the view that schools occupy relatively
persistent positions in the distribution of structural characteristics, and
support the use of culture scores measured in 2024--25 as proxies for the
organisational conditions present during the 2023--24 outcome period.

The same analysis at the level of the thirteen primary-need categories
within the SEN population shows moderate stability for the Social, Emotional
and Mental Health (SEMH) category (ICC~$= 0.545$), which supports the
baseline SEMH share entering as an additional control in the robustness
specification of \cref{tab:robustness_overall}. The lower stability of SEMH
relative to the aggregate SEN support share (ICC~$= 0.65$) is consistent
with SEMH composition being partly responsive to school behaviour and
exclusion practices, rather than being determined solely by fixed intake
characteristics.

\begin{table}[htb]

\centering\small
\begin{tabular}{lrrrr}
\toprule
Variable & \multicolumn{1}{c}{Schools} & \multicolumn{1}{c}{Mean yrs}  & \multicolumn{1}{c}{ICC} & \multicolumn{1}{c}{Mean YoY} \\
\midrule
EAL share (\%)  & {3,286} & 7.7 & 0.922 & 3.7 pp \\
Pupils on roll  & {3,363} & 9.8 & 0.920 & 37 pupils \\
KS2 scaled score (2021--24)  & {3,215} & 3.0 & 0.910 & 0.86 pts \\
KS2 APS (2015--19)  & {3,100} & 4.0 & 0.812 & 0.65 pts \\
FSM share (\%)  & {3,363} & 8.8 & 0.742 & 2.3 pp \\
Progress~8  & {3,256} & 6.8 & 0.727 & 0.21 pts \\
EHC plan share (\%)  & {3,355} & 9.6 & 0.669 & 0.4 pp \\
SEN support share (\%)  & {3,361} & 9.8 & 0.653 & 1.9 pp \\
Per-pupil expenditure  & {3,331} & 6.7 & 0.516 & \pounds599 \\
\bottomrule
\end{tabular}
\caption{Structural stability of school characteristics}
\label{tab:structural_stability}
\begin{minipage}{\linewidth}
\smallskip
\footnotesize\textit{Notes:} The panel covers 2015--16 to 2024--25. ICC = intra-class correlation = between-school
variance / (between-school + within-school variance), estimated via one-way
random-effects ANOVA on the school-year panel; schools with fewer than three
years of non-missing data are excluded. \textit{Mean yrs}: mean number of
academic years of non-missing data per school. \textit{Mean YoY}: mean
absolute year-on-year change per school. EAL = English as an Additional
Language; FSM = free school meals (school census); EHC = Education, Health and
Care plan; SEN = special educational needs. The two KS2 measures use different scales from successive assessment frameworks and are therefore reported separately; no KS4 performance data
were published for 2019--20 or 2020--21 owing to the suspension of national
examinations. Progress~8 is nationally standardised with a mean near zero;
the mean YoY figure is an absolute change, not a percentage.
Per-pupil expenditure combines Consistent Financial Reporting (maintained
schools) and Academies Accounts Return (academies); the 2015--16 and 2024--25 expenditure data are incomplete.
\end{minipage}
\end{table}

\Cref{tab:p8_component_stability} reports ICC estimates separately for the four
Progress~8 components that serve as the primary outcomes of this chapter.
Among the Progress~8 components,
stability varies from ICC~$= 0.587$ for the Open component to ICC~$= 0.748$
for EBaC, with English and Maths sitting between these at ICC~$\approx 0.67$.
The higher stability of EBaC relative to Open reflects the fact that
curriculum-pathway decisions, that is, which subjects pupils are entered for,
are governed by stable school policy and change little from year to year,
while the more varied Open bucket (arts, vocational and additional GCSEs) is
subject to greater annual variation in subject mix. All four P8 components
exceed the per-pupil expenditure ICC of~$0.52$, which is the weakest result
in the main stability table; the cross-sectional design is therefore at least
as defensible for each component individually as for the structural
characteristics already discussed.

The theoretical prediction of this chapter is that warmth is most strongly
associated with the Open component, which is the dimension most sensitive to
breadth of provision and pupil engagement, and that strictness is most
strongly associated with the EBaC component, which is the dimension most
sensitive to academic pathway steering. The lower stability of Open P8 relative to EBaC P8
(ICC~$= 0.587$ versus $0.748$) implies that the warmth coefficient will be
estimated against a noisier outcome, which may attenuate the estimated effect.
This is a limitation to acknowledge in interpreting any differential finding
between components.

\begin{table}[htb]
\centering\small
\begin{tabular}{lrrrr}
\toprule
Variable & \multicolumn{1}{c}{Schools} & \multicolumn{1}{c}{Mean yrs} & \multicolumn{1}{c}{ICC} & \multicolumn{1}{c}{Mean YoY} \\
\midrule
\quad English   & {3,256} & 6.8 & 0.666 & 0.250 pts \\
\quad Maths     & {3,256} & 6.8 & 0.667 & 0.225 pts \\
\quad EBacc     & {3,256} & 6.8 & 0.748 & 0.243 pts \\
\quad Open      & {3,256} & 6.8 & 0.587 & 0.304 pts \\
\quad Overall   & {3,256} & 6.8 & 0.727 & 0.210 pts \\
\bottomrule
\end{tabular}
\caption{Structural stability of the Progress~8 components}
\label{tab:p8_component_stability}
\begin{minipage}{\linewidth}
\smallskip
\footnotesize\textit{Notes:} ICC and mean years are defined as in
\cref{tab:structural_stability}. The panel covers 2015--16 to 2023--24,
excluding 2019--20 and 2020--21, for which no KS4 performance data were
published owing to the suspension of national examinations. The components
are English (two slots, double-weighted), Maths (two slots,
double-weighted), EBacc (sciences, computer science, history, geography, and
modern foreign languages), and Open (any additional approved
qualification).
\end{minipage}
\end{table}


% - - - - - - - - - - - - - - - - - - - - - - - - - - - - - - -
\section{Additional Tables}
\label{app:3A:tables}

\Cref{tab:main_results_s1,tab:main_results_s2,tab:main_results_s3} report the
full five-outcome versions of the three main-body stages, including the EBaC
and Open components. These tables use the primary specification ($n = 99$,
predecessor-filled pre-COVID grades, late-entry schools excluded), which is
the same one as the main-body tables, so the Overall, English and Maths
columns reproduce the body exactly and the two extra columns can be read on
the same footing.
Reading the EBaC and Open columns, warmth and strictness are both
significant for EBaC with near-equal magnitudes, and both are significant for
Open as well, again with near-equal magnitudes. As \cref{app:3A:p8mechanics} explains, both components mix progress with
the school's entry decisions, so the entry-rate analysis in
\cref{tab:entry_rates} is the clearer way to answer the pathway question.
The strictness score predicts EBacc, humanities and languages entry
nationally with the inspection grade held constant, but controlling for
EBacc entry attenuates the strictness--EBaC association by only 11 to
14~per~cent. This means that most of that association is progress within
subjects rather than curriculum steering.

\begin{table}[htbp]\centering
\def\sym#1{\ifmmode^{#1}\else\(^{#1}\)\fi}
\caption{Stage 1: Total culture effect ($W + S$). Primary specification: full control set,
predecessor-filled pre-COVID inspection grade, late-entry schools excluded ---
identical to the trio reported in the chapter body, extended to all five
Progress~8 components.}
\label{tab:main_results_s1}
\resizebox{\textwidth}{!}{%
\begin{tabular}{l*{5}{c}}
\toprule
                    &     Overall         &     English         &       Maths         &        EBaC         &        Open         \\
\midrule
Warmth ($W$)        &       0.143\sym{***}&       0.167\sym{***}&       0.112\sym{**} &       0.160\sym{***}&       0.136\sym{**} \\
                    &     (0.045)         &     (0.043)         &     (0.054)         &     (0.057)         &     (0.061)         \\
\addlinespace
Strictness ($S$)    &       0.131\sym{***}&       0.087\sym{*}  &       0.110\sym{**} &       0.169\sym{***}&       0.139\sym{**} \\
                    &     (0.047)         &     (0.049)         &     (0.051)         &     (0.058)         &     (0.066)         \\
\midrule
$N$                 &          99         &          99         &          99         &          99         &          99         \\
$R^2$               &       0.624         &       0.621         &       0.527         &       0.596         &       0.518         \\
\bottomrule
\multicolumn{6}{l}{\footnotesize Standard errors in parentheses}\\
\multicolumn{6}{l}{\footnotesize \sym{*} \(p<0.10\), \sym{**} \(p<0.05\), \sym{***} \(p<0.01\)}\\
\end{tabular}%
}
\end{table}

\begin{table}[htbp]\centering
\def\sym#1{\ifmmode^{#1}\else\(^{#1}\)\fi}
\caption{Stage 2: Teaching quality benchmark ($T$)}
\label{tab:main_results_s2}
\resizebox{\textwidth}{!}{%
\begin{tabular}{l*{5}{c}}
\toprule
                    &     Overall         &     English         &       Maths         &        EBaC         &        Open         \\
\midrule
Teaching quality ($T$)&       0.251\sym{***}&       0.232\sym{***}&       0.205\sym{***}&       0.281\sym{***}&       0.274\sym{***}\\
                    &     (0.052)         &     (0.053)         &     (0.060)         &     (0.066)         &     (0.058)         \\
\midrule
$N$                 &          96         &          96         &          96         &          96         &          96         \\
$R^2$               &       0.579         &       0.571         &       0.489         &       0.522         &       0.527         \\
\bottomrule
\multicolumn{6}{l}{\footnotesize Standard errors in parentheses}\\
\multicolumn{6}{l}{\footnotesize \sym{*} \(p<0.10\), \sym{**} \(p<0.05\), \sym{***} \(p<0.01\)}\\
\end{tabular}%
}
\end{table}

\begin{table}[htbp]\centering
\def\sym#1{\ifmmode^{#1}\else\(^{#1}\)\fi}
\caption{Stage 3: Culture net of teaching ($W + S + T$). Primary specification: full control set,
predecessor-filled pre-COVID inspection grade, late-entry schools excluded ---
identical to the trio reported in the chapter body, extended to all five
Progress~8 components.}
\label{tab:main_results_s3}
\resizebox{\textwidth}{!}{%
\begin{tabular}{l*{5}{c}}
\toprule
                    &     Overall         &     English         &       Maths         &        EBaC         &        Open         \\
\midrule
Warmth ($W$)        &       0.130\sym{**} &       0.172\sym{***}&       0.103         &       0.170\sym{**} &       0.085         \\
                    &     (0.061)         &     (0.063)         &     (0.065)         &     (0.077)         &     (0.075)         \\
\addlinespace
Strictness ($S$)    &       0.110\sym{***}&       0.078\sym{*}  &       0.092\sym{*}  &       0.152\sym{***}&       0.102\sym{*}  \\
                    &     (0.041)         &     (0.044)         &     (0.049)         &     (0.052)         &     (0.057)         \\
\addlinespace
Teaching quality ($T$)&       0.020         &      -0.011         &       0.014         &      -0.021         &       0.086         \\
                    &     (0.064)         &     (0.068)         &     (0.073)         &     (0.078)         &     (0.083)         \\
\midrule
$N$                 &          99         &          99         &          99         &          99         &          99         \\
$R^2$               &       0.624         &       0.621         &       0.527         &       0.596         &       0.526         \\
\bottomrule
\multicolumn{6}{l}{\footnotesize Standard errors in parentheses}\\
\multicolumn{6}{l}{\footnotesize \sym{*} \(p<0.10\), \sym{**} \(p<0.05\), \sym{***} \(p<0.01\)}\\
\end{tabular}%
}
\end{table}

\begin{table}[htbp]\centering
\def\sym#1{\ifmmode^{#1}\else\(^{#1}\)\fi}
\small
\footnotesize\setlength{\tabcolsep}{4pt}
\begin{tabular}{llcccc}
\toprule
\multicolumn{6}{l}{\textit{Panel A: entry-rate regressions}} \\
\addlinespace[2pt]
 & Predictor & $N$ & EBacc entry & Humanities entry & Languages entry \\
\midrule
A & Visited warmth ($W$) & 100 & 0.001 (0.033) & $-$0.008 (0.014) & 0.017 (0.030) \\
A & Visited strictness ($S$) & 100 & 0.036 (0.032) & 0.026\sym{**} (0.012) & 0.014 (0.030) \\
A & Marking-scheme bands & 3{,}076 & 0.033\sym{***} (0.004) & 0.016\sym{***} (0.002) & 0.029\sym{***} (0.004) \\
\addlinespace
B & Visited warmth ($W$) & 99 & $-$0.006 (0.033) & $-$0.009 (0.015) & 0.010 (0.030) \\
B & Visited strictness ($S$) & 99 & 0.045 (0.033) & 0.027\sym{**} (0.013) & 0.021 (0.031) \\
B & Marking-scheme bands & 2{,}753 & 0.022\sym{***} (0.004) & 0.014\sym{***} (0.003) & 0.018\sym{***} (0.004) \\
\midrule
\multicolumn{6}{l}{\textit{Panel B: channel decomposition, strictness $\rightarrow$ EBacc Progress 8 component}} \\
\addlinespace[2pt]
 & Model & $N$ & Without entry & With entry & Attenuation \\
\midrule
A & National (marking-scheme bands) & 3{,}052 & 0.134\sym{***} & 0.115\sym{***} & 14.4\% \\
A & Visited ($S$) & 100 & 0.232\sym{***} & 0.205\sym{***} & 11.3\% \\
B & National (marking-scheme bands) & 2{,}746 & 0.112\sym{***} & 0.100\sym{***} & 10.2\% \\
B & Visited ($S$) & 99 & 0.244\sym{***} & 0.218\sym{***} & 10.7\% \\
\bottomrule
\end{tabular}
\caption{Culture and curriculum entry}
\label{tab:entry_rates}
\begin{minipage}{\linewidth}
\smallskip
\footnotesize\textit{Notes:} Panel A asks whether warm or strict schools enter pupils for more academic curricula: each cell regresses an entry rate (EBacc, humanities, languages) on the standardised culture score plus controls, for the visit scores and the national marking-scheme strictness bands, without (A) and with (B) the predecessor-filled 2019 inspection-grade control; late-entry schools excluded. Panel B then asks how much of the strictness--EBacc-progress association runs through entry: the strictness coefficient on the EBacc Progress 8 component before and after controlling the EBacc entry rate. Standard errors (HC3) in parentheses.
\sym{*} \(p<0.10\), \sym{**} \(p<0.05\), \sym{***} \(p<0.01\).
The national strictness instrument is the marking-scheme band of Chapter~2
(1--5, standardised over each estimation sample).
Entry rates are derived from pupil counts. 
\end{minipage}
\end{table}


\Cref{tab:outcome_stats_full} reports descriptive statistics for all five
Progress~8 components and the four Attainment~8 buckets, and
\cref{tab:national_strictness_full} the national extension on all five
outcomes; the EBaC and Open coefficients are of the same size as the others.

\begin{table}[htbp]
  \centering
\caption{Progress~8 and Attainment~8 outcomes: summary statistics}
  \label{tab:outcome_stats_full}
  \footnotesize\setlength{\tabcolsep}{4pt}
  \begin{tabular}{llrrrrr}
    \toprule
    Outcome & Sample & $N$ & Mean & SD & Min & Max \\
    \midrule
    \multirow{3}{*}{P8 Overall (2-yr avg)} & Visited & 103 & 0.252 & 0.478 & -0.700 & 1.345 \\
     & Interviewed & 299 & 0.095 & 0.487 & -1.655 & 1.345 \\
     & National & 3240 & -0.015 & 0.510 & -3.570 & 2.465 \\
    \addlinespace[3pt]
    \multirow{3}{*}{P8 English (2-yr avg)} & Visited & 103 & 0.256 & 0.507 & -0.750 & 1.545 \\
     & Interviewed & 299 & 0.093 & 0.485 & -1.485 & 1.545 \\
     & National & 3240 & -0.017 & 0.519 & -4.035 & 2.360 \\
    \addlinespace[3pt]
    \multirow{3}{*}{P8 Maths (2-yr avg)} & Visited & 103 & 0.270 & 0.449 & -0.595 & 1.345 \\
     & Interviewed & 299 & 0.093 & 0.470 & -1.370 & 1.345 \\
     & National & 3240 & -0.017 & 0.487 & -2.030 & 2.660 \\
    \addlinespace[3pt]
    \multirow{3}{*}{P8 EBacc (2-yr avg)} & Visited & 103 & 0.297 & 0.549 & -0.790 & 1.490 \\
     & Interviewed & 299 & 0.116 & 0.570 & -1.960 & 1.490 \\
     & National & 3240 & -0.020 & 0.587 & -4.950 & 2.655 \\
    \addlinespace[3pt]
    \multirow{3}{*}{P8 Open (2-yr avg)} & Visited & 103 & 0.172 & 0.528 & -1.095 & 1.435 \\
     & Interviewed & 299 & 0.055 & 0.554 & -3.095 & 1.500 \\
     & National & 3240 & -0.034 & 0.573 & -5.140 & 2.220 \\
    \addlinespace[3pt]
    \multirow{3}{*}{Att8 English (2024--25)} & Visited & 103 & 10.658 & 1.757 & 7.900 & 17.300 \\
     & Interviewed & 300 & 10.277 & 1.696 & 6.700 & 17.300 \\
     & National & 3274 & 10.049 & 1.702 & 0.100 & 17.300 \\
    \addlinespace[3pt]
    \multirow{3}{*}{Att8 Maths (2024--25)} & Visited & 103 & 9.887 & 1.975 & 6.700 & 17.400 \\
     & Interviewed & 300 & 9.514 & 1.891 & 5.400 & 17.400 \\
     & National & 3274 & 9.263 & 1.879 & 0.300 & 17.600 \\
    \addlinespace[3pt]
    \multirow{3}{*}{Att8 EBacc (2024--25)} & Visited & 103 & 14.783 & 3.235 & 9.400 & 26.500 \\
     & Interviewed & 300 & 14.156 & 3.118 & 6.200 & 26.500 \\
     & National & 3274 & 13.696 & 3.103 & 0.500 & 26.600 \\
    \addlinespace[3pt]
    \multirow{3}{*}{Att8 Open (2024--25)} & Visited & 103 & 14.676 & 2.936 & 10.000 & 26.300 \\
     & Interviewed & 300 & 14.184 & 2.868 & 4.900 & 26.300 \\
     & National & 3274 & 13.869 & 2.830 & 0.000 & 26.300 \\
    \bottomrule
  \end{tabular}
  \begin{minipage}{\linewidth}
    \vspace{2pt}
    \footnotesize
    \textit{Notes:} The interviewed sample numbers 303 schools and contains the 103 visited schools; $N$ falls short of 303 where an outcome is unpublished for a school. P8 outcomes are 2-year averages of 2022--23 and 2023--24 component scores where both years are available. Att8 2024--25 components are the
    contemporaneous robustness outcome; no Progress~8 is available for this cohort because the cohort were
    in Year~6 in 2019--20 when Key Stage~2 assessments were cancelled.  
  \end{minipage}
\end{table}

% generated by analyse_band_extras.py -- do not edit
\begin{table}[htbp]
\centering
\caption{The national strictness leg (marking-scheme bands) on all five Progress 8 outcomes, full controls without the grade, HC3.}
\label{tab:national_strictness_full}
\begin{tabular}{lccc}
\toprule
Outcome & Coefficient & (se) & $n$ \\
\midrule
Overall & +0.130*** & (0.008) & 3,052 \\
English & +0.114*** & (0.008) & 3,052 \\
Maths & +0.110*** & (0.008) & 3,052 \\
EBaC & +0.135*** & (0.009) & 3,052 \\
Open & +0.153*** & (0.009) & 3,052 \\
\bottomrule
\end{tabular}
\end{table}


\Cref{tab:subscores} decomposes the visit-based culture scores into their
five observational blocks and enters them jointly. Only two survive:
outside-lesson warmth and in-lesson strictness. The progress association
runs through relational quality in corridors and transitions, and orderly
conduct in classrooms, rather than through their counterparts.
\Cref{tab:items_fdr} takes the decomposition to the level of individual
observation items, each entered in its own small model with a
false-discovery correction across the set and each item's inter-rater
reliability shown beside it. Every one of the 33 items has a positive
coefficient and 28 survive the false-discovery correction, so the sub-score
result reflects which facets contain \emph{independent} signal, and not
which items relate to progress at all.

\begin{table}[htbp]\centering
\def\sym#1{\ifmmode^{#1}\else\(^{#1}\)\fi}
\small
\caption{The five visit sub-scores entered jointly. Classroom warmth (W1), outside-of-classroom
warmth (W2), classroom strictness (S1), outside-of-classroom strictness (S2) and teaching
quality (T1) are entered together, each standardised within the estimation sample, with the
primary control set. Outside-of-lesson warmth and in-lesson strictness carry the association;
the other three add little once those two are held fixed.}
\label{tab:subscores}
\begin{tabular}{lccc}
\toprule
Sub-score & Overall & English & Maths \\
\midrule
Classroom warmth (W1)        & 0.061 & 0.105 & 0.112 \\
                             & (0.080) & (0.078) & (0.084) \\
\addlinespace
Outside warmth (W2)          & 0.108\sym{**} & 0.144\sym{***} & 0.075 \\
                             & (0.047) & (0.047) & (0.051) \\
\addlinespace
Classroom strictness (S1)    & 0.120\sym{***} & 0.137\sym{***} & 0.105\sym{*} \\
                             & (0.044) & (0.041) & (0.054) \\
\addlinespace
Outside strictness (S2)      & 0.005 & $-$0.053 & $-$0.004 \\
                             & (0.045) & (0.046) & (0.052) \\
\addlinespace
Teaching quality (T1)        & 0.006 & $-$0.052 & $-$0.047 \\
                             & (0.082) & (0.088) & (0.087) \\
\midrule
$N$                          & 99 & 99 & 99 \\
\bottomrule
\end{tabular}
\begin{minipage}{\linewidth}
\smallskip
\footnotesize\textit{Notes:} Standard errors (HC3) in parentheses.
\sym{*} \(p<0.10\), \sym{**} \(p<0.05\), \sym{***} \(p<0.01\).
Primary specification: visited schools, late-entry schools excluded, full control set with
the predecessor-filled 2019 grade. The sub-scores are correlated, so the individual
coefficients should be read with care: variance inflation factors are 5.79 (W1),
1.93 (W2), 1.96 (S1), 2.08 (S2) and 6.55 (T1) --- W1 and T1
overlap heavily, which widens their standard errors.
\end{minipage}
\end{table}

\begin{table}[htbp]\centering
\def\sym#1{\ifmmode^{#1}\else\(^{#1}\)\fi}
\footnotesize
\caption{Each visit-sheet item as a predictor of Progress 8, one small model per item.
Every row regresses average Progress 8 on that item alone plus the full control set with
the predecessor-filled grade control ($N=100$). Rows are sorted by coefficient size.
The $q$ column is the Benjamini--Hochberg false-discovery-rate adjusted $p$-value across
all 33 items; a dagger marks items that survive at $q<0.05$. The final columns give each
item's inter-rater reliability: the pooled pairwise Pearson correlation between
independent observers, and the number of observer pairs it rests on.}
\label{tab:items_fdr}
\begin{tabular}{llcccccr}
\toprule
Item & Setting & Sub-score & $\beta$ & SE & $q$ & IRR $r$ & Pairs \\
\midrule
Outcomes        & Classroom & T1 & 0.213\sym{\dagger} & 0.038 & $<$0.001 & 0.466 & 468 \\
Concentration   & Classroom & S1 & 0.196\sym{\dagger} & 0.036 & $<$0.001 & 0.647 & 467 \\
Motivation      & Classroom & W1 & 0.183\sym{\dagger} & 0.042 & $<$0.001 & 0.610 & 467 \\
Explanation     & Classroom & T1 & 0.170\sym{\dagger} & 0.039 & $<$0.001 & 0.532 & 467 \\
Respectfully    & Classroom & S1 & 0.169\sym{\dagger} & 0.044 & $<$0.001 & 0.654 & 467 \\
Students        & Outside   & W2 & 0.166\sym{\dagger} & 0.038 & $<$0.001 & 0.705 & 93  \\
Response        & Classroom & S1 & 0.159\sym{\dagger} & 0.038 & $<$0.001 & 0.668 & 468 \\
Misbehaviour    & Classroom & S1 & 0.158\sym{\dagger} & 0.035 & $<$0.001 & 0.693 & 469 \\
Methods         & Classroom & T1 & 0.158\sym{\dagger} & 0.041 & $<$0.001 & 0.572 & 468 \\
Questioning     & Classroom & T1 & 0.157\sym{\dagger} & 0.045 & 0.001    & 0.585 & 468 \\
Relationships   & Outside   & W2 & 0.144\sym{\dagger} & 0.035 & $<$0.001 & 0.444 & 93  \\
Disruption      & Classroom & S1 & 0.142\sym{\dagger} & 0.041 & 0.001    & 0.789 & 468 \\
Verbal          & Classroom & T1 & 0.140\sym{\dagger} & 0.051 & 0.009    & 0.523 & 468 \\
Praise          & Classroom & W1 & 0.140\sym{\dagger} & 0.042 & 0.002    & 0.656 & 466 \\
Structure       & Classroom & T1 & 0.139\sym{\dagger} & 0.038 & $<$0.001 & 0.443 & 468 \\
Differentiation & Classroom & T1 & 0.137\sym{\dagger} & 0.049 & 0.008    & 0.543 & 468 \\
Names           & Classroom & W1 & 0.135\sym{\dagger} & 0.037 & $<$0.001 & 0.351 & 466 \\
Student         & Classroom & W1 & 0.135\sym{\dagger} & 0.047 & 0.006    & 0.591 & 466 \\
Resource        & Classroom & T1 & 0.132\sym{\dagger} & 0.042 & 0.003    & 0.437 & 468 \\
Corridors       & Outside   & S2 & 0.130\sym{\dagger} & 0.045 & 0.006    & 0.657 & 93  \\
Sanction        & Outside   & S2 & 0.128\sym{\dagger} & 0.043 & 0.005    & 0.678 & 93  \\
Arrival         & Outside   & S2 & 0.122\sym{\dagger} & 0.041 & 0.005    & 0.399 & 93  \\
Teacher         & Classroom & T1 & 0.122\sym{\dagger} & 0.044 & 0.008    & 0.617 & 468 \\
Interactions    & Outside   & W2 & 0.120\sym{\dagger} & 0.046 & 0.012    & 0.476 & 93  \\
Interactions    & Classroom & W1 & 0.118\sym{\dagger} & 0.043 & 0.009    & 0.567 & 467 \\
Reward          & Outside   & -- & 0.105\sym{\dagger} & 0.040 & 0.011    & 0.738 & 93  \\
Verbal          & Outside   & -- & 0.096           & 0.051 & 0.068    & 0.398 & 90  \\
Displays        & Outside   & -- & 0.090\sym{\dagger} & 0.037 & 0.017    & 0.489 & 93  \\
Recreational    & Outside   & S2 & 0.087\sym{\dagger} & 0.039 & 0.032    & 0.486 & 93  \\
Differentiation & Outside   & -- & 0.080           & 0.052 & 0.128    & 0.360 & 89  \\
Canteen         & Outside   & S2 & 0.056           & 0.036 & 0.128    & 0.551 & 93  \\
Discussion      & Classroom & T1 & 0.053           & 0.045 & 0.247    & 0.765 & 467 \\
Alignment       & Outside   & -- & 0.042           & 0.041 & 0.307    & 0.320 & 93  \\
\bottomrule
\end{tabular}
\begin{minipage}{\linewidth}
\smallskip
\footnotesize\textit{Notes:} \sym{\dagger} survives Benjamini--Hochberg correction at
$q<0.05$ (28 of 33 items). Items marked ``--'' in the sub-score column are recorded on
the visit sheet but do not enter any composite. Gold-standard visited tier, late-entry
schools excluded. The composites remain the headline measures; these per-item models are
mechanism texture, not competing specifications.
\end{minipage}
\end{table}


\Cref{tab:typology} reports the combination analysis. In the visited sample
the warmth-by-strictness interaction is null while the warm-and-strict
quadrant tops the adjusted means, so the two dimensions add rather than
amplify; the national panel using the inspection-report scores shows the
same additive pattern once the grade is controlled. \Cref{tab:gaps} examines
whether the gap between a school's espoused and enacted culture predicts
progress; it does not, and the one apparently significant gap term is shown
in the final columns to be an artefact of document-score extremity rather
than of mismatch itself.

\begin{table}[htbp]\centering
\def\sym#1{\ifmmode^{#1}\else\(^{#1}\)\fi}
\small
\footnotesize\setlength{\tabcolsep}{4pt}
\begin{tabular}{lccc}
\toprule
 & Visited schools & \multicolumn{2}{c}{National (text instruments)} \\
\cmidrule(lr){3-4}
 & & Panel A & Panel B \\
\midrule
\multicolumn{4}{l}{\textit{Adjusted quadrant means (mean, SE, $n$)}} \\
\addlinespace[2pt]
Authoritative (high $W$, high $S$) & $+$0.547 (0.066) & $+$0.347 (0.025) & $+$0.275 (0.026) \\
   & $n=33$ & $n=299$ & $n=277$ \\
Authoritarian (low $W$, high $S$)  & $+$0.142 (0.075) & $+$0.240 (0.047) & $+$0.213 (0.049) \\
   & $n=16$ & $n=46$ & $n=44$ \\
Permissive (high $W$, low $S$)     & $+$0.380 (0.073) & $+$0.022 (0.010) & $+$0.053 (0.011) \\
   & $n=15$ & $n=1{,}227$ & $n=1{,}095$ \\
Neglectful (low $W$, low $S$)      & $+$0.035 (0.058) & $-$0.078 (0.009) & $-$0.051 (0.009) \\
   & $n=35$ & $n=1{,}480$ & $n=1{,}330$ \\
\midrule
\multicolumn{4}{l}{\textit{Regression tests}} \\
\addlinespace[2pt]
Warmth ($W$)             & 0.150\sym{***} & 0.054\sym{***} & 0.055\sym{***} \\
                         & (0.043) & (0.007) & (0.008) \\
Strictness ($S$)         & 0.109\sym{**} & 0.117\sym{***} & 0.096\sym{***} \\
                         & (0.042) & (0.008) & (0.008) \\
$W \times S$ interaction & 0.039 & 0.004 & $-$0.005 \\
                         & (0.035) & (0.008) & (0.008) \\
Authoritative vs rest    & 0.403\sym{***} & 0.363\sym{***} & 0.262\sym{***} \\
                         & (0.088) & (0.027) & (0.027) \\
\midrule
$N$                      & 99 & 3{,}052 & 2{,}746 \\
\bottomrule
\end{tabular}
\caption{The authoritative-school typology tested directly. The upper panel gives
covariate-adjusted mean Progress 8 for the four quadrants formed by splits on
warmth and strictness; the lower panel gives the continuous interaction test and the
contrast of the authoritative (high-warmth, high-strictness) quadrant against all
others. The first column uses the visit scores on the primary
specification; the national columns use the prediction-model warmth and the
marking-scheme strictness bands of Chapter~2, without (A) and with (B) the
predecessor-filled 2019 inspection-grade control. In the national columns warmth is
split at the national median of the predicted score, and ``high'' strictness means
above the in-sample band median of 4 --- that is, band 5 only --- so the
high-strictness cells are small. Scores are standardised within each estimation
sample; late-entry schools excluded throughout.}
\label{tab:typology}
\begin{minipage}{\linewidth}
\smallskip
\footnotesize\textit{Notes:} Standard errors (HC3) in parentheses; for the
visited-sample adjusted quadrant means they are HC1, because the HC3 leverage
adjustment is undefined when a covariate cell contains a single school.
\sym{*} \(p<0.10\), \sym{**} \(p<0.05\), \sym{***} \(p<0.01\).
Quadrant means are covariate-adjusted; visited-school quadrants are median splits within the
visited sample. The visited-sample interaction is a low-power test at this sample size; the
pattern --- interaction near zero with the high-both quadrant on top --- reads as
additive rather than synergistic, and the national columns repeat it: the
interaction is indistinguishable from zero while the high-both quadrant sits on
top. The national columns use the prediction-model warmth and the marking-scheme
strictness bands of Chapter~2, not the retired analyser scores; because ``high''
strictness is band 5 only, the two high-strictness cells are thin and their means
carry wider standard errors.
\end{minipage}
\end{table}

\begin{table}[htbp]\centering
\caption{The espoused--enacted gap and Progress~8}
\label{tab:gaps}
\def\sym#1{\ifmmode^{#1}\else\(^{#1}\)\fi}
\small

\footnotesize\setlength{\tabcolsep}{4pt}
\begin{tabular}{lccc}
\toprule
 & Signed gap & Absolute gap & \shortstack{Absolute gap\\(+ component quadratics)} \\
\midrule
\multicolumn{4}{l}{\textit{Visited sample (visit vs interview, $n=99$)}} \\
\addlinespace[2pt]
Warmth & 0.047 & 0.032 & 0.011 \\
 & (0.040) & (0.046) & (0.058) \\
Strictness & 0.065 & 0.002 & $-$0.015 \\
 & (0.050) & (0.061) & (0.061) \\
\bottomrule
\end{tabular}
\begin{minipage}{\linewidth}
\smallskip
\footnotesize\textit{Notes:} Each row regresses average Progress 8 on the enacted (observed) score, the espoused-minus-enacted gap, and the primary control set, on the visited schools. The first column enters the signed gap, the second the absolute gap, and the third re-enters the absolute gap alongside squared terms in both the espoused and the enacted score. An absolute gap could predict progress simply because progress is a non-linear function of one of the scores. The squared terms remove that explanation. Standard errors in parentheses (HC3).
\sym{*} \(p<0.10\), \sym{**} \(p<0.05\), \sym{***} \(p<0.01\).
Scores are z-standardised before the gap is formed; late-entry schools excluded. 
\end{minipage}
\end{table}


\Cref{tab:llm_p8_matrix} reports, for completeness, the association of every
text-derived instrument with the five outcomes under the grade-controlled
specification. Only the inspection-report scores have substantial, robust
associations. The website binary indicators have small ones, and the
behaviour policy warmth score is significantly negative, which is the
policy-as-aspiration finding of \cref{ch:paper1} reappearing as a
predictive result. \Cref{tab:parentview} gives the parent-survey panel.

\begin{table}[htbp]\centering
\def\sym#1{\ifmmode^{#1}\else\(^{#1}\)\fi}
\caption{Every LLM-derived culture score against every Progress 8 outcome. Each cell is
the coefficient from one regression of the outcome on the (standardised) predictor plus
the national control set, including the 2019 inspection-grade control. Because the sweep runs 150 such
tests (75 per grade panel), significance is marked only after Benjamini--Hochberg
correction across the whole matrix. The Ofsted-derived scores carry the grade-reconstruction
caution described in the text; behaviour-policy rows use standard errors clustered on
the shared document digest.}
\label{tab:llm_p8_matrix}
\resizebox{\textwidth}{!}{%
\begin{tabular}{lrccccc}
\toprule
Predictor & $N$ & Overall & English & Maths & EBacc & Open \\
\midrule
Ofsted strictness          & 2{,}823 & 0.122\sym{\dagger} & 0.110\sym{\dagger} & 0.109\sym{\dagger} & 0.125\sym{\dagger} & 0.141\sym{\dagger} \\
                           &         & (0.007) & (0.008) & (0.008) & (0.009) & (0.008) \\
Ofsted warmth              & 2{,}823 & 0.073\sym{\dagger} & 0.062\sym{\dagger} & 0.055\sym{\dagger} & 0.073\sym{\dagger} & 0.093\sym{\dagger} \\
                           &         & (0.006) & (0.007) & (0.006) & (0.007) & (0.008) \\
Ofsted teaching            & 2{,}823 & 0.118\sym{\dagger} & 0.109\sym{\dagger} & 0.103\sym{\dagger} & 0.123\sym{\dagger} & 0.131\sym{\dagger} \\
                           &         & (0.007) & (0.008) & (0.007) & (0.009) & (0.009) \\
BP strictness (v4)         & 2{,}822 & $-$0.003 & $-$0.001 & 0.001 & $-$0.005 & $-$0.005 \\
                           &         & (0.007) & (0.007) & (0.007) & (0.008) & (0.008) \\
BP warmth (v4)             & 2{,}822 & $-$0.016\sym{\dagger} & $-$0.009 & $-$0.011 & $-$0.014 & $-$0.026\sym{\dagger} \\
                           &         & (0.007) & (0.007) & (0.007) & (0.008) & (0.008) \\
Website warmth (v13)       & 2{,}832 & $-$0.001 & 0.001 & $-$0.005 & 0.002 & $-$0.003 \\
                           &         & (0.006) & (0.007) & (0.006) & (0.007) & (0.007) \\
Website strictness (v13)   & 2{,}832 & 0.003 & $-$0.000 & 0.010 & 0.002 & 0.002 \\
                           &         & (0.007) & (0.007) & (0.007) & (0.008) & (0.009) \\
Website strictness (v15)   & 2{,}832 & 0.003 & $-$0.000 & 0.009 & 0.001 & 0.002 \\
                           &         & (0.007) & (0.007) & (0.007) & (0.008) & (0.009) \\
Traditional ethos (v2)     & 2{,}832 & 0.028\sym{\dagger} & 0.022\sym{\dagger} & 0.032\sym{\dagger} & 0.029\sym{\dagger} & 0.029\sym{\dagger} \\
                           &         & (0.007) & (0.008) & (0.008) & (0.009) & (0.008) \\
Traditional pedagogy (v1b) & 2{,}832 & 0.028\sym{\dagger} & 0.028\sym{\dagger} & 0.025\sym{\dagger} & 0.036\sym{\dagger} & 0.025\sym{\dagger} \\
                           &         & (0.009) & (0.009) & (0.009) & (0.010) & (0.009) \\
Faith prominence (0--3)    & 2{,}832 & 0.027\sym{\dagger} & 0.029\sym{\dagger} & $-$0.004 & 0.014 & 0.063\sym{\dagger} \\
                           &         & (0.007) & (0.008) & (0.007) & (0.008) & (0.009) \\
Interview strictness (v13) & 254     & 0.003 & 0.001 & $-$0.002 & 0.021 & $-$0.009 \\
                           &         & (0.026) & (0.026) & (0.027) & (0.030) & (0.029) \\
Interview teaching (v3)    & 254     & 0.019 & 0.005 & 0.030 & 0.040 & $-$0.002 \\
                           &         & (0.026) & (0.026) & (0.026) & (0.029) & (0.035) \\
Interview warmth (v15)     & 254     & 0.027 & 0.027 & 0.010 & 0.044 & 0.020 \\
                           &         & (0.024) & (0.024) & (0.026) & (0.029) & (0.029) \\
Interview warmth (counts)  & 254     & 0.041 & 0.044 & $-$0.005 & 0.040 & 0.071\sym{\dagger} \\
                           &         & (0.022) & (0.023) & (0.025) & (0.025) & (0.030) \\
\bottomrule
\end{tabular}%
}
\begin{minipage}{\linewidth}
\smallskip
\footnotesize\textit{Notes:} Standard errors in parentheses (HC3 throughout, except the
behaviour-policy rows, which are clustered on the shared trust-template document
digest). \sym{\dagger} survives Benjamini--Hochberg correction at $q<0.05$ across all 150
tests in the sweep (both grade panels); the panel shown here includes the 2019
inspection-grade control, which is why $N$ falls below the full spine. The companion
panel without the grade control gives the same qualitative picture with slightly larger
Ofsted coefficients. Interview rows are limited to schools with a scoreable
head-teacher interview.
\end{minipage}
\end{table}

\begin{table}[htbp]\centering
\def\sym#1{\ifmmode^{#1}\else\(^{#1}\)\fi}
\small
\caption{Parent View as a soft national criterion for the warmth instruments. The parent
composite is the mean of the ``my child is happy at this school'' and ``my child feels
safe at this school'' positive-response rates. The upper panel correlates each warmth
instrument with the composite and then regresses the composite on the instrument plus
the control set; the lower panel reports the composite's own association with Progress 8,
which is descriptive only because parents observe outcomes.}
\label{tab:parentview}
\begin{tabular}{lccccc}
\toprule
\multicolumn{6}{l}{\textit{Panel A: warmth instruments against the parent composite}} \\
\addlinespace[2pt]
Instrument & $r$ & $\rho$ & $n$ & Controlled $\beta$ & $p$ \\
\midrule
Website warmth (v18)          & $+$0.024 & $+$0.024 & 3{,}235 & $+$0.001 (0.003) & 0.71 \\
Website warmth (v13)          & $+$0.024 & $+$0.021 & 3{,}235 & $+$0.001 (0.003) & 0.78 \\
Behaviour-policy warmth (v4)  & $-$0.038 & $-$0.041 & 3{,}225 & $-$0.005 (0.003) & 0.07 \\
Ofsted warmth                 & $+$0.401 & $+$0.395 & 3{,}212 & $+$0.052 (0.003) & $<$0.001 \\
\midrule
\multicolumn{6}{l}{\textit{Panel B: parent composite as a Progress 8 predictor (descriptive)}} \\
\addlinespace[2pt]
Outcome & \multicolumn{2}{c}{$\beta$} & SE & $n$ & $p$ \\
\midrule
P8 overall  & \multicolumn{2}{c}{$+$0.137} & 0.008 & 2{,}750 & $<$0.001 \\
P8 English  & \multicolumn{2}{c}{$+$0.109} & 0.009 & 2{,}750 & $<$0.001 \\
P8 Maths    & \multicolumn{2}{c}{$+$0.114} & 0.008 & 2{,}750 & $<$0.001 \\
\bottomrule
\end{tabular}
\begin{minipage}{\linewidth}
\smallskip
\footnotesize\textit{Notes:} Standard errors in parentheses. Panel A regressions control
for the full national control set. Parent View is a soft criterion: it conflates school
culture with parent satisfaction, each school contributes its single most recent survey
release (releases span September 2022 to September 2025), and responses are not
temporally aligned with the instrument sources. The document-based instruments
(website, behaviour policy) do not track parents' experienced climate; the Ofsted-derived
score does, consistent with both drawing on visits to the school. Panel B is endogenous
--- parents plausibly respond to results as well as climate --- so it is reported as a
description, never as a causal predictor.
\end{minipage}
\end{table}


\Cref{tab:p8_proxy} reports the 2024--25 pseudo-progress robustness check:
an imputed value-added outcome for the one cohort measured contemporaneously
with the fieldwork, validated by running the identical construction one year
earlier, where it correlates with the published Progress~8 at $+0.81$. The
culture coefficients on the pseudo-outcome sit within $0.05$ of their
real-P8 counterparts. \Cref{tab:stability_p8} supports the outcome-year
strategy directly: within-school Progress~8 movement across years is small,
and the academisation event-time columns show why conversion effects must
be read against conversion type. Sponsor-led schools decline before
converting and recover afterwards, which is the pattern one would expect
from regression to the mean, while converters are flat throughout.

\begin{table}[htbp]\centering
\def\sym#1{\ifmmode^{#1}\else\(^{#1}\)\fi}
\small
\footnotesize\setlength{\tabcolsep}{4pt}
\begin{tabular}{llccccc}
\toprule
 & Predictor & $N$ & \shortstack{Pseudo\\overall} & \shortstack{Pseudo\\English} & \shortstack{Pseudo\\Maths} & \shortstack{Real P8\\overall} \\
\midrule
A & Visited warmth ($W$)   & 100     & 0.162\sym{***} & 0.182\sym{***} & 0.114\sym{*} & 0.154\sym{***} \\
  &                        &         & (0.050) & (0.051) & (0.061) & (0.041) \\
A & Visited strictness ($S$) & 100     & 0.135\sym{***} & 0.092\sym{*} & 0.156\sym{***} & 0.098\sym{**} \\
  &                        &         & (0.050) & (0.050) & (0.056) & (0.040) \\
A & Ofsted strictness      & 3{,}076 & 0.154\sym{***} & 0.134\sym{***} & 0.136\sym{***} & 0.137\sym{***} \\
  &                        &         & (0.009) & (0.009) & (0.009) & (0.007) \\
\addlinespace
B & Visited warmth ($W$)   & 99      & 0.165\sym{***} & 0.195\sym{***} & 0.116\sym{*} & 0.147\sym{***} \\
  &                        &         & (0.057) & (0.057) & (0.067) & (0.043) \\
B & Visited strictness ($S$) & 99      & 0.135\sym{**} & 0.081 & 0.157\sym{**} & 0.107\sym{**} \\
  &                        &         & (0.055) & (0.055) & (0.061) & (0.041) \\
B & Ofsted strictness      & 2{,}753 & 0.153\sym{***} & 0.136\sym{***} & 0.136\sym{***} & 0.136\sym{***} \\
  &                        &         & (0.009) & (0.010) & (0.010) & (0.007) \\
\bottomrule
\end{tabular}
\caption{Culture coefficients on the 2024/25 stand-in outcome}
\label{tab:p8_proxy}
\begin{minipage}{\linewidth}
\smallskip
\footnotesize\textit{Notes:} The pseudo measure residualises 2024/25 Attainment 8 on a multi-year z-averaged KS2 intake proxy, giving a progress-like score for the cohort measured contemporaneously with the fieldwork. Each cell is the (standardised) culture coefficient from the same specification run on the pseudo outcome and on real (2022--24 average) Progress 8 for the same schools, without (A) and with (B) the predecessor-filled 2019 inspection-grade control; late-entry schools excluded. Standard errors (HC3) in parentheses.
\sym{*} \(p<0.10\), \sym{**} \(p<0.05\), \sym{***} \(p<0.01\).
Across the national panel the pseudo measure correlates with real average Progress 8 at
$r = +0.719$ ($n = 3{,}140$). Validation of the construction one year back (the
pseudo measure rebuilt from 2021--22 and 2022--23 baselines against the real 2023--24
Progress 8) is $+0.81$ (two-year construction, $n = 3{,}242$;
ch3\_appendix\_p8proxy\_validation.csv). On the body's per-standard-deviation scale, the
visited-school coefficients on the pseudo outcome sit within $0.04$ of their
real-Progress 8 twins in the same panel (Panel A: pseudo 0.162/0.135 against real
0.154/0.098; Panel B: pseudo 0.165/0.135 against real 0.147/0.107,
warmth/strictness). The ``Ofsted strictness'' rows use the marking-scheme strictness
bands of Chapter~2 (1--5, standardised over each estimation sample), not the retired
analyser score. The pseudo measure is
school-level, not pupil-matched, so it is a robustness check, never a primary outcome.
\end{minipage}
\end{table}

\begin{table}[htbp]\centering
\small

\begin{tabular}{lrrr}
\toprule
\multicolumn{4}{l}{\textit{Panel A: within-school variation (2017--18 to 2023--24, $n=3{,}232$ schools)}} \\
\addlinespace[2pt]
\multicolumn{4}{l}{Within-school SD: mean 0.196, median 0.177, 10th percentile 0.086, 90th percentile 0.330} \\
\addlinespace[4pt]
Adjacent-year pair & $N$ & Moved $>$0.3 & Moved $>$0.5 \\
\midrule
2017--18 $\rightarrow$ 2018--19 & 2{,}762 & 19.0\% & 4.5\% \\
2021--22 $\rightarrow$ 2022--23 & 3{,}185 & 21.9\% & 5.1\% \\
2022--23 $\rightarrow$ 2023--24 & 3{,}126 & 18.5\% & 3.1\% \\
All adjacent pairs pooled       & 9{,}073 & 19.9\% & 4.3\% \\
\midrule
\multicolumn{4}{l}{\textit{Panel B: mean Progress 8 by event year around academisation}} \\
\addlinespace[2pt]
Event year & \multicolumn{1}{c}{Converter} & \multicolumn{1}{c}{Sponsor-led} & \\
\midrule
$t-5$ & $-$0.08 \,($n=135$) & $-$0.42 \,($n=71$) & \\
$t-4$ & $-$0.08 \,($n=121$) & $-$0.45 \,($n=68$) & \\
$t-3$ & $-$0.11 \,($n=165$) & $-$0.52 \,($n=123$) & \\
$t-2$ & $-$0.12 \,($n=246$) & $-$0.54 \,($n=197$) & \\
$t-1$ & $-$0.10 \,($n=335$) & $-$0.53 \,($n=242$) & \\
$t+0$ & $-$0.08 \,($n=359$) & $-$0.43 \,($n=274$) & \\
$t+1$ & $-$0.03 \,($n=293$) & $-$0.38 \,($n=226$) & \\
$t+2$ & $-$0.07 \,($n=207$) & $-$0.33 \,($n=165$) & \\
$t+3$ & $-$0.08 \,($n=167$) & $-$0.32 \,($n=137$) & \\
$t+4$ & $-$0.06 \,($n=145$) & $-$0.27 \,($n=136$) & \\
$t+5$ & $-$0.05 \,($n=195$) & $-$0.27 \,($n=181$) & \\
$t+6$ & $-$0.03 \,($n=196$) & $-$0.29 \,($n=175$) & \\
$t+7$ & $-$0.03 \,($n=132$) & $-$0.31 \,($n=107$) & \\
$t+8$ & $+$0.01 \,($n=56$) & $-$0.21 \,($n=43$) & \\
\bottomrule
\end{tabular}
\caption{Stability of Progress~8 and academisation}
\label{tab:stability_p8}
\begin{minipage}{\linewidth}
\smallskip
\footnotesize\textit{Notes:} Panel A summarises each school's own standard deviation of Progress 8 across the five published years and the share of schools moving by more than 0.3 or 0.5 of a grade between consecutive years. Panel B gives mean Progress 8 by year relative to the year of academy conversion, separately for converter and sponsor-led academies. Panel A uses schools with at least three published years.
There is no Progress 8 for 2019--20 or 2020--21 (COVID), so ``adjacent'' means
consecutive academic years only; the 2018--19 to 2021--22 boundary is not counted as a
pair. Panel B rests on the predecessor--successor bridging table (converter and
sponsor-led academies; cells with fewer than five schools suppressed); these are raw
event-time means, not a matched event study. 
\end{minipage}
\end{table}


\Cref{tab:semh_mechanism} tests the sorting concern directly, regressing the
2023--24 share of pupils with social, emotional and mental-health needs on
strictness, conditional on the school's 2015--16 share and the standard
controls. Nationally, stricter schools have a lower current share conditional
on their baseline ($-0.108$ of a percentage point per standard deviation of
the score, $p = 0.013$, $n = 2{,}594$); the direction is the same in the visited
sample but imprecise ($-0.257$, $p = 0.36$, $n = 94$), and
warmth is unrelated to the change in either sample. The pattern is consistent
with sorting, though it cannot distinguish selective admission from schools
supporting pupils to leave the register. Adding the baseline share to the
primary specification (\cref{tab:robustness_overall}) moves the
strictness coefficient from $0.106$ to $0.090$; re-estimating on the same 94
schools without the control gives $0.097$, so about seven per cent of the
strictness association is attributable to the control itself and the rest of
the movement to the change of sample.

\begin{table}[htbp]\centering
\caption{The SEMH composition mechanism}
\label{tab:semh_mechanism}
\def\sym#1{\ifmmode^{#1}\else\(^{#1}\)\fi}

\begin{tabular}{l*{3}{c}}
\toprule
                    &\multicolumn{1}{c}{Visited (S)}&\multicolumn{1}{c}{Visited (W)}&\multicolumn{1}{c}{National (S)}\\
\midrule
Strictness ($S_{\text{visit}}$, per SD)&      $-$0.257         &                     &                     \\
                    &     (0.280)         &                     &                     \\
\addlinespace
SEMH share 2015--16 (\%)&       0.176 &       0.191\sym{*}&       0.194\sym{***}\\
                    &     (0.120)         &     (0.115)         &     (0.022)         \\
\addlinespace
Warmth ($W_{\text{visit}}$, per SD)&                     &      $-$0.094         &                     \\
                    &                     &     (0.240)         &                     \\
\addlinespace
Strictness (marking-scheme bands, per SD)&                     &                     &      $-$0.114\sym{***}\\
                    &                     &                     &     (0.038)         \\
\midrule
N                   &          94         &          94         &        2{,}594         \\
$R^2$                  &       0.559         &       0.547         &       0.484         \\
\bottomrule
\end{tabular}
\begin{minipage}{\linewidth}
\smallskip
\footnotesize\textit{Notes:} Each column regresses a school's 2023--24 share of pupils with social, emotional and mental-health (SEMH) needs, as a percentage of the roll, on one culture score plus the school's 2015--16 baseline SEMH share and the primary control set: visited-school strictness and warmth ($n=94$ visited schools), and the national marking-scheme strictness bands ($n=2{,}594$). Coefficients are percentage points per standard deviation of the score. The strictness coefficients are null or slightly negative: there is a detectable trace of sorting, and controlling for it removes under a tenth of the strictness--progress association.
Standard errors (HC3) in parentheses. Outcome: 2023--24 SEMH pupils as a percentage of the roll; baseline is the 2015--16 SEMH count over the current roll. Primary controls; late-entry excluded. \sym{*} \(p<0.10\), \sym{**} \(p<0.05\), \sym{***} \(p<0.01\).
\end{minipage}
\end{table}


\Cref{tab:score_controls_corr} reports pairwise Pearson correlations between
the three enacted culture scores ($W_{\text{enac}}$,
$S_{\text{enac}}$, $T_{\text{enac}}$) and the two espoused scores
($W_{\text{esp}}$, $S_{\text{esp}}$) against each control variable,
estimated on the Tier~1 sample of \NTierOne{} visited schools. No pairwise
correlation exceeds $|r| = 0.29$. The largest entries are between
$W_{\text{enac}}$ and EAL share ($r = -0.25$), between $T_{\text{enac}}$ and
EAL share ($r = -0.26$) and KS2 prior attainment ($r = +0.24$), and between
$W_{\text{enac}}$ and selective school status ($r = +0.20$). The absence of
high correlations indicates that the culture scores are not collinear with
the observable school characteristics included in the regression. Selective
schools score modestly higher on both culture dimensions ($r = +0.20$ for
warmth and $+0.22$ for strictness), which the selective-admissions control
absorbs in every specification.

\begin{table}[htbp]
  \centering
  \caption{Pairwise correlations: culture scores and control variables (Tier~1, $N=103$)}
  \label{tab:score_controls_corr}
  \small
  \begin{tabular}{lrrrrr}
    \toprule
    Control & $W_{\text{enac}}$ & $S_{\text{enac}}$ & $T_{\text{enac}}$ & $W_{\text{esp}}$ & $S_{\text{esp}}$ \\
    \midrule
    KS2 & +0.19 & +0.15 & +0.24 & +0.15 & $-$0.15 \\
    FSM\% & $-$0.22 & $-$0.04 & $-$0.03 & $-$0.11 & +0.12 \\
    EAL\% & $-$0.29 & +0.00 & $-$0.26 & $-$0.04 & +0.03 \\
    SEN\% & $-$0.10 & $-$0.11 & $-$0.04 & $-$0.12 & $-$0.08 \\
    Log size & +0.01 & $-$0.04 & +0.06 & $-$0.10 & $-$0.06 \\
    Academy & +0.04 & $-$0.07 & +0.06 & $-$0.08 & +0.05 \\
    Urban & $-$0.20 & $-$0.04 & $-$0.16 & $-$0.10 & $-$0.05 \\
    Selective & $-$0.23 & $-$0.07 & $-$0.07 & $-$0.06 & $-$0.09 \\
    Yrs Ofsted & $-$0.03 & $-$0.09 & $-$0.12 & $-$0.07 & $-$0.04 \\
    \bottomrule
  \end{tabular}
  \begin{minipage}{\linewidth}
    \vspace{2pt}
    \footnotesize
    \textit{Notes:} Pearson correlation coefficients. $W_{\text{enac}}$,
    $S_{\text{enac}}$ and $T_{\text{enac}}$ are the enacted culture scores
    built from school visit observations only (mean of $W_1, W_2$; $S_1, S_2$;
    $T_1$, each scaled 0--10). $W_{\text{esp}}$ and $S_{\text{esp}}$ are the
    espoused scores built from the headteacher interview. The two sets share no
    component; they are reported side by side rather than averaged, for the
    reasons given in \cref{sec:p1_composite}.
    Bold entries indicate $|r| \geq 0.30$. $N$ varies slightly by column due to missing
    control values; all correlations use the Tier~1 sample of 103 visited schools.
  \end{minipage}
\end{table}


A final check re-estimates the Stage~1 specification on the subset of
schools where the headteacher named in the most recent inspection report is
confirmed still in post in the July~2026 GIAS snapshot. Continuity could be
determined for only 29 of the 103 full-data schools, of which 22 were
unchanged and 7 had changed, so the restriction is driven as much by data
coverage as by turnover. In this subsample the standard errors are roughly
four times larger than in the full sample, and neither warmth nor strictness
is statistically significant in any component. The strictness point
estimates are in the same direction as the full-sample estimates and are
larger in every component, while the warmth estimates are noisy and should
not be interpreted directionally. This check does not provide evidence
against the primary results, but the subsample is too small to detect
effects of the size found in the full specification.
